% Modernised reading — fonds Alexandre Grothendieck, shelfmark 36, pages 1-84
% (the whole folder, five batches, all transcribed).
%
% This is an INTERPRETATION, not a transcription. It re-reads the pages in
% today's notation and vocabulary, states the hypotheses the manuscript leaves
% implicit, and drops the critical apparatus so that the mathematics reads
% straight through. Everything the brackets and underlines carried moves into a
% footnote. It is held to being mathematically correct as it stands; where the
% pages are loose or mistaken, the true statement is what appears here and the
% footnote says what the page has. In any dispute of substance the
% transcription governs: whoever cites Grothendieck must cite batch-NN.fr.tex.
%
% Pass: Opus 5.5 (claude-opus-5-5), 2026-10-04 — first pass, unchecked against the pages by a human.
% Checked: Opus 5.5 (claude-opus-5-5), 2026-10-09 — the skill's checks run, the reading re-read against
% the transcriptions, five local corrections (pp. 82, 70, 78 footnotes; spine; p. 27).
%
% Scope: batches 1-5, pages 1-84, all transcribed, read whole before any of
% this was written; ONE file for the shelfmark.
%
% What the folder is. His working papers for his share of SGA 7 (Exposés VII-IX:
% biextensions, Néron models and monodromy), plus one typescript that is not
% part of any published exposé. Seven runs, most under a cover in his hand:
% (1) pp. 2-7, an unidentified reader's remarks on a typed state of Exposé IX,
%     with his pencil replies — another person's text, noted only, as the
%     transcription notes it (RIGHTS.md);
% (2) pp. 9-26, the typescript « Le théorème de Griffiths par voie
%     algébrique », §§ 1-7 and a bare heading 8;
% (3) pp. 27-33, notes on biextensions of complexes [M -> G] (what Deligne later
%     named 1-motives), on the backs of three typed leaves of Exposé VIII;
% (4) pp. 34-45, « Picard et dualité »;
% (5) pp. 46-57, « Biextensions de faisceaux de groupes »;
% (6) pp. 59-70, « Th. du cycle invariant : cas du H^1 », typescript and
%     manuscript working;
% (7) pp. 72-80, « Th. d'orthogonalité pour les modèles de Néron,
%     démonstration arithmétique »;
% (8) pp. 81-83, « Autodualité de la jacobienne, suivant Artin-Mazur ».
% Leaves of other people's typescripts (Bourbaki pp. 37, 65-71; a Čech
% cohomology typescript pp. 48, 50; drafts of Mumford's GIT pp. 52, 54, 56, 58)
% are mentioned only as the transcription mentions them.
%
% Notation fixed for the whole document (departures from the pages' letters):
% — In the Griffiths typescript, sections 4-7: dim X = m throughout. The page
%   uses m in § 4, 2n in §§ 5-6 and n in § 7 (where its n is § 4's m); the
%   reading writes m everywhere and m = 2n where §§ 5-6 need it. In 7.6 the sum
%   of degrees the page calls r (also the ambient P^r) is e; in the proof of 7.1
%   q = p^f (the page: p^m). The structural morphism of the curve C in § 1,
%   which the page leaves unnamed and later calls g, is q_C : C -> S; g is
%   X -> S.
% — Pages 35-45: his curly z_l, z_p are written mu_l, mu_p and his i_p is
%   alpha_p — the values his table gives are exactly theirs; his f_!, f^!, f_.
%   in the adjunction formulas are written Rf_* (the right adjoint of f^*).
% — Pages 73-76: the subgroup the page calls S (in G_Y) is H, to free S.
% — Nothing else is renamed.
%
% Substantive departures from the pages, each footnoted where it occurs:
% — p. 9: the degree condition m = m' + N - 2d' (the typed « - dim(Y'/S) » is
%   struck and left blank).
% — p. 13: Y_s -> X_s for the page's X_s -> X_s.
% — p. 14: « Grif(Y) (x) Q contains P^{2n}_alg(X) as a subspace » is stated as
%   what the argument gives: a subspace of Grif (x) Q mapping ONTO P^{2n}_alg(X).
% — p. 15: E^{2n-1} for the page's E^{n-1}.
% — p. 17: the generic hypersurface's coefficients are indexed by monomials of
%   degree d, not t_0..t_r.
% — pp. 17-26: Théorème 7.1 is given in the form the ink amends it to (trace on
%   a u-stable subspace V'_0); the typed form (trace on all of V_0), which 7.2
%   needs, is not established by the argument, and his own margins say
%   « pas prouvé » of 7.1 and 7.2. Lemme 7.8's « sous-groupe ouvert » is a
%   decomposition group (closed, not open in general).
% — p. 23: Lemme 7.6 needs N <= e d_1^m; the page has the inequality reversed
%   (its own remark « N <= r » for d_1 = 1 confirms the direction).
% — p. 30: n-divisibility nG = G for the page's « _nG = G ».
% — p. 38: the inductive limit of M/nM is M (x) Q/Z, not M (x) Q; and the two
%   limit sequences give H^1(X, A) = the TORSION subgroup of Ext^1(A', Pic_X).
% — p. 55: the quadratic identity is phi(u+v) + phi(u-v) = 2(phi(u)+phi(v)).
% — p. 57: k o Delta^* j = 2 splits the extension only after pull-back by 2.
% — p. 60: d^t = p^h d with h >= 0, and p^h divisible by pgcd(mu_i) — not equal
%   to it (counterexample in the footnote).
% — p. 62: Ker(Br(K) -> Br(X_eta)) for the page's X_etabar.
% — p. 63: d(X'/S') is the p-part of d (= p^r only when d^t = d).
%
% What the folder announces and never establishes: §§ 7-8 of the Griffiths
% typescript beyond what is above (7.2 and 7.3 rest on the unproved typed 7.1;
% § 8 is a heading); the proofs of a)-d) and Corollaires 9, 10, 15 of the
% invariant-cycle typescript (the manuscript proves the case n prime to p and
% refers b) to Raynaud); the equivalences of Corollaires 1-2 of p. 47-49; the
% questions of pp. 39, 40 and 70 (« Sous quelles conditions ??? », « demander à
% Raynaud si Ker est tout Z/dZ »).
\documentclass[11pt,a4paper]{article}
\input{../preamble/grothendieck}

\folder{36}
\pages{1}{84}
\dating{[vers 1967-1973]}
\watermark{Édition de démonstration\\interprétation personnelle de l'œuvre}
\foldertitle{SGA 7 (ma part) : notes manuscrites (s.d.), tapuscrit annoté (s.d.) — lecture modernisée du dossier entier}

\begin{document}

\section*{Résumé}

\begin{resume}
Ces feuillets sont les papiers de travail de Grothendieck pour sa part d'un
séminaire collectif, SGA~7, consacré à ce qui arrive à une famille de variétés
algébriques quand l'un de ses membres dégénère.

L'image la plus simple est celle d'une famille de tores --- de surfaces en forme
de bouée --- dont l'un, au centre, est pincé : un de ses cercles s'est réduit à
un point. Sur un tore voisin, deux sortes de lacets cohabitent. Celui qui va se
pincer, le \emph{cycle évanescent}, revient à lui-même quand on fait tourner le
paramètre autour du centre ; l'autre, qui fait le tour dans l'autre sens, revient
déformé, alourdi d'une copie du premier. Cette transformation est la
\emph{monodromie}. Ce que le dossier cherche à comprendre, c'est l'algèbre de
cette situation quand le tore est remplacé par une variété abélienne ou une
jacobienne, et les lacets par un module de Tate : quels lacets sont fixes, quels
lacets « meurent » dans la fibre pincée, et comment la dualité --- l'accouplement
qui mesure l'enlacement de deux lacets --- se comporte vis-à-vis de ces deux
sous-ensembles. La réponse tient en un mot : \emph{orthogonalité}. Les lacets qui
meurent d'un côté sont orthogonaux aux lacets fixes de l'autre ; dans l'exemple
du tore, le lacet qui se pince n'enlace pas lui-même.

Pour le dire, il faut un outil, et le dossier le forge en même temps : la
\emph{biextension}, une façon d'encoder un accouplement bilinéaire entre deux
groupes qui garde la trace de ce qui se passe au-dessus de la fibre pincée. Une
bonne part des notes en développe le maniement : comment une biextension se
reconstruit à partir de ses morceaux, comment on la décrit quand les groupes sont
des complexes plutôt que des groupes, comment la forme d'intersection des
composantes d'une fibre dégénérée produit, sur le groupe des composantes, une
dualité finie qui mesure exactement l'obstruction à prolonger la dualité
au-dessus du point singulier.

Un texte tapé tranche sur le reste : un essai pour démontrer « par voie
algébrique », en toute caractéristique, un théorème de Griffiths --- qu'une
variété peut porter des cycles qui sont des bords sans pouvoir être déformés en
rien. L'idée est de compter des points sur des corps finis : un nombre de points
qui n'est pas ce qu'il « devrait » être modulo $p$ trahit une cohomologie qui
n'est pas de la forme qu'on voulait exclure.

Ce dossier montre aussi un mathématicien qui relit son propre texte et le
démonte. Sur le tapuscrit où il écrit « je vais prouver », l'encre ajoute
« non, pas tout à fait ! », puis, en face du théorème et du corollaire, « pas
prouvé ». Ailleurs : « la triste vérité est que je ne sais pas montrer » ; « Sous
quelles conditions ??? » ; « demander à Raynaud ». Ce qu'il appelait l'attitude
d'écoute passe aussi par là : laisser le texte dire où il ne tient pas, plutôt
que de le faire tenir.

Les noms modernes à chercher ensuite sont ceux de la dernière ligne.

\keywords{Griffiths group, coniveau, Hodge level, Lefschetz pencil, vanishing cycles, Woods Hole fixed point formula, Chebotarev density theorem, 1-motive, biextension, Weil pairing, fppf cohomology, Cartier duality, Artin-Schreier-Witt theory, Néron-Severi group, invariant cycle theorem, Néron model, Grothendieck's orthogonality theorem, Frobenius weights, component group, Grothendieck's pairing}
\end{resume}

\section*{Le fil du dossier, et les conventions}

Le titre de l'inventaire --- « SGA 7 (ma part) » --- dit ce que le dossier
rassemble : ce dont Grothendieck s'était chargé dans le séminaire sur la
monodromie, tenu à Bures en 1967--1969 et publié en deux volumes (SGA~7~I, 1972 ;
SGA~7~II, 1973). Sa part publiée comprend les exposés I, II et VI à IX ; ceux que
touche ce dossier sont VII et VIII (biextensions) et IX (modèles de Néron et
monodromie). Le fil qui relie les feuillets est le
suivant : une variété abélienne $A$ sur le corps des fractions $K$ d'un anneau de
valuation discrète, son module de Tate $M = T_\ell(A)$, la filtration
\[ M \supset M^{I} \supset M^{I}_{t} \]
par la partie fixe sous l'inertie et la partie torique, et la façon dont
l'accouplement de Weil entre $M$ et le module de la variété duale traite ces
filtrations. Tout le reste est outillage, ou détour.

Les stations, dans l'ordre des pages :
\begin{itemize}
\item \textbf{Pages 2 à 7} --- les remarques d'un lecteur sur un état tapé de
l'exposé IX, avec ses réponses au crayon. Un texte d'autrui, signalé seulement.
\item \textbf{Pages 9 à 26} --- le tapuscrit « Le théorème de Griffiths par voie
algébrique », n\textsuperscript{os} 1 à 7 et un titre de n\textsuperscript{o}~8.
C'est le détour : il relève de l'autre moitié de SGA~7 (pinceaux de Lefschetz,
cycles évanescents) et d'aucun exposé publié.
\item \textbf{Pages 27 à 33} --- biextensions de complexes $[M \to G]$, l'objet
que Deligne nommera plus tard 1-motif.
\item \textbf{Pages 34 à 45} --- « Picard et dualité » : cohomologie plate d'une
courbe, torseurs sous un groupe fini et sous une variété abélienne, groupe
fondamental d'une variété abélienne.
\item \textbf{Pages 46 à 57} --- « Biextensions de faisceaux de groupes » :
reconstruction d'une biextension par restriction, biextensions sans
rigidification, le groupe de Néron--Severi.
\item \textbf{Pages 59 à 70} --- le théorème des cycles invariants pour le
$H^1$, en tapuscrit puis à la main.
\item \textbf{Pages 72 à 80} --- le théorème d'orthogonalité, démonstration
arithmétique.
\item \textbf{Pages 81 à 83} --- l'autodualité du groupe des composantes du
modèle de Néron d'une jacobienne.
\end{itemize}
Les feuillets d'autres textes que le dossier contient --- une rédaction Bourbaki
(page 37 et pages 65 à 71), un tapuscrit sur la cohomologie de Čech (pages 48 et
50), des feuillets d'une rédaction de la théorie géométrique des invariants de
Mumford (pages 52, 54, 56, 58) --- ne sont signalés qu'à leur place, et pas
au-delà de ce qu'en dit la transcription.

\emph{Conventions.} $k$ désigne le corps de base ou le corps résiduel selon la
section, ce que chaque section rappelle ; $\ell$ est un nombre premier
différent de la caractéristique résiduelle quand rien d'autre n'est dit. Pour un
groupe abélien (ou un faisceau) $P$ et un entier $n$, on écrit comme lui
${}_nP$ le noyau et $P_n = P/nP$ le conoyau de la multiplication par $n$. Dans
le tapuscrit Griffiths, la dimension de la variété ambiante $X$ est notée $m$
partout\footnote{Le tapuscrit la note $m$ au n\textsuperscript{o}~4, $2n$ aux
n\textsuperscript{os}~5 et 6, et $n$ au n\textsuperscript{o}~7, où son $n$ est le
$m$ du n\textsuperscript{o}~4. Garder ses lettres aurait fait désigner par $n$
deux entiers différents à deux pages d'intervalle.} ; on écrit $m = 2n$ là où les
n\textsuperscript{os}~5 et 6 en ont besoin. Les faisceaux qu'il note d'un $z$
bouclé (pages 35 et suivantes) sont écrits $\mu_\ell$, $\mu_p$, et son $i_p$
est écrit $\alpha_p$ : les valeurs que sa table leur donne sont exactement les
leurs.

\pagerange{2}{7}
\section*{Les remarques d'un lecteur sur l'exposé IX (pages 2 à 7)}

Six feuillets quadrillés, numérotés 1 à 6, titrés « Remarques. SGA~7 IX », d'une
main non identifiée qui tutoie l'auteur de l'exposé : un relecteur parcourt un
état dactylographié de l'exposé IX --- \emph{Modèles de Néron et monodromie},
l'exposé de Grothendieck --- section par section, en renvoyant aux pages et aux
lignes du tapuscrit. Ce texte n'est pas le sien et n'est pas reproduit ici ; il
suffit de dire qu'il couvre l'introduction et les sections 1 à 11 (la 8 laissée
vide), et qu'au feuillet 2 le lecteur compare l'accouplement « d'A.~M. » à celui de
l'auteur.\footnote{Que « A.~M. » soit Artin--Mazur, dont le dernier titre du
dossier (page 81) se réclame, est notre inférence ; ni les remarques ni la
transcription ne le disent.}

Ce qui est de lui, ce sont les réponses au crayon, et elles disent assez bien
comment l'exposé a été défendu : « guillemets suffisants ? », « je le dis ! »,
« c'est dit dans la 5.13 », « c'est pourtant trivial », « non, mon argument
marche », « si ? », deux fois « que faire ? », et, à l'encre, en face d'une
objection qui jugeait fausse une page entière, « CANULAR ?! ».\footnote{Le crayon
du feuillet 2 a été vérifié de sa main sur le fac-similé ; pour les autres
feuillets la transcription le dit vraisemblablement le sien, sans plus, et l'encre
de « CANULAR ?! » « présumément » sienne.}

\pagerange{9}{26}
\section*{Le théorème de Griffiths par voie algébrique (pages 9 à 26)}

Un tapuscrit de dix-huit pages, paginé à la main, sous un titre au crayon. Il est
écrit à la première personne par un auteur qui n'est ni Deligne ni Katz, il parle
de « la théorie des cycles évanescents, qui sera prouvée dans SGA 7 cette
année », et ce qu'il démontre n'est dans aucun des exposés publiés de
Grothendieck.\footnote{Le plan de SGA~7 du dossier 35 portait « XX Th.\ de
Griffiths », biffé. Dans SGA~7~II publié, l'exposé XX, de Katz, s'intitule, à
notre connaissance, « Le niveau de la cohomologie des intersections complètes »,
et l'exposé XXII, du même, « Une formule de congruence pour la fonction $\zeta$ » :
ce sont les deux ingrédients que ce tapuscrit appelle « le procédé de Katz » et
« la congruence de Deligne ». Le rapprochement est nôtre. La phrase sur SGA~7
« cette année » situe le tapuscrit pendant le séminaire, 1967--1969.} Le
problème est celui que Griffiths avait résolu par voie transcendante en 1969 :
exhiber des cycles homologiquement équivalents à zéro qui ne sont pas
algébriquement équivalents à zéro. Le tapuscrit veut le faire en toute
caractéristique.

Le plan est simple à énoncer. Une classe de cohomologie primitive, sur l'espace
total d'une famille, donne une classe de $H^1$ de la base (n\textsuperscript{o}~1).
Si un cycle de la fibre générique est algébriquement trivial, cette classe vient de
la variété ambiante, sous deux hypothèses a) et b) (n\textsuperscript{os}~2 et 3).
Pour un pinceau de Lefschetz, la classe d'un cycle primitif ne vient jamais de la
variété ambiante (n\textsuperscript{o}~4). Donc, sous a) et b), un cycle primitif
non nul de $X$ se restreint en un cycle non trivial du groupe de Griffiths de la
section hyperplane générique (n\textsuperscript{o}~5). Reste à vérifier b), ce qui
est l'objet d'une conjecture (n\textsuperscript{o}~6) et d'un théorème annoncé
(n\textsuperscript{o}~7) que l'auteur, à l'encre, déclare non prouvé.

\pagerange{9}{11}
\subsection*{La classe $u(x)$ d'une classe primitive (n\textsuperscript{o}~1, pages 9 à 11)}

Soit $f\colon Y \to S$ un morphisme propre et lisse, $S$ régulier et connexe, à
caractéristiques résiduelles différentes de $\ell$, et $F$ un faisceau
$\ell$-adique lisse de rang un sur $S$ (en pratique $\mathbf{Z}_\ell(i)$). Une
classe $x \in H^m(Y, F_Y)$ est dite \emph{primitive} si son image dans
$H^0(S, R^m f_* F)$ est nulle ; comme $R^m f_* F$ est lisse et $S$ connexe, il
revient au même que sa restriction à une fibre géométrique $Y_{\bar s}$ soit
nulle. Dans la suite spectrale de Leray
$E_2^{p,q} = H^p(S, R^q f_*F) \Rightarrow H^{p+q}(Y, F)$, une classe primitive est
dans le premier cran $F^1H^m$ de la filtration, et son image dans
$F^1/F^2 = E_\infty^{1,m-1}$ est un élément
\[ u(x) \in E_2^{1,m-1} = H^1(S, R^{m-1}f_*F_Y), \]
puisqu'aucune différentielle n'aboutit en bidegré $(1, m-1)$ et que
$E_\infty^{1,m-1}$ est donc un sous-groupe de $E_2^{1,m-1}$.\footnote{Le
tapuscrit pose $u(x) \in E_2^{1,m-1}$ sans ce commentaire. L'exposant $m$ est
partout écrit ou repassé à la main sur une autre lettre tapée.} Appliquée à la
classe d'un cycle, c'est l'avatar $\ell$-adique de la \emph{fonction normale}
de Poincaré et Griffiths ; aujourd'hui on parle d'application d'Abel--Jacobi
$\ell$-adique.

\emph{Fonctorialité.} Soit $Y'$ propre et lisse sur $S$, de dimension relative
constante $d'$, et $Z \in H^N(Y \times_S Y', G)$, avec un accouplement
$F' \otimes G \to F(d')$. Par le formalisme des correspondances cohomologiques,
$Z$ définit $x' \mapsto Z(x') = \mathrm{pr}_{1*}(\mathrm{pr}_2^*x' \cdot Z)$, de
$H^{m'}(Y', F')$ dans $H^m(Y, F)$, avec
\[ m = m' + N - 2d' . \]
\footnote{Le tapuscrit écrit « $m = m' + N - \dim(Y'/S)$ », puis la machine
biffe « $\dim(Y'/S)$ » sans rien mettre à la place. L'image directe le long d'une
fibration de dimension relative $d'$ abaisse le degré de $2d'$ : c'est ce terme qui
manque, et la torsion $(d')$ de l'accouplement, que le tapuscrit donne, en est la
contrepartie.} Si $x'$ est primitive, $Z(x')$ l'est, et
\[ u(Z(x')) = Z(u(x')) , \]
où le second membre se lit à travers l'homomorphisme
$R^{m'-1}f'_*F_{Y'} \to R^{m-1}f_*F_Y$ que $Z$ définit par localisation étale sur
$S$.

\emph{Le cas qui servira.} On prend pour $Y'$ une courbe relative
$q_C\colon C \to S$ à fibres géométriques connexes, munie de deux sections $t_0$,
$t_1$, et pour $x' \in H^2(C, \mathbf{Z}_\ell(1))$ la classe du diviseur relatif
$t_1(S) - t_0(S)$, primitive parce que de degré nul sur chaque fibre ; on prend
pour $Z$ la classe d'un cycle de codimension $n$ sur $Y \times_S C$, encore noté
$Z$.\footnote{Le tapuscrit ne nomme pas le morphisme structural de $C$, et
l'écrit $g$ à la page 11 ; on le note $q_C$, la lettre $g$ servant plus loin à
$X \to S$.} Alors $x = Z(x')$ est la classe du cycle $z = Z(t_1) - Z(t_0)$ de
$Y$, de degré $2n$, et $u(x) = Z(u(x'))$ dans $H^1(S, R^{2n-1}f_*\mathbf{Z}_\ell(n))$.
Le tapuscrit prend soin de dire que ces classes de cycles et ces intersections
n'ont de sens, « en l'état actuel de la Science », que sous des hypothèses du
genre « $S$ de type fini sur un corps », et on les suppose. La conséquence qui
servira est celle-ci :

\begin{quote}
si l'homomorphisme $Z\colon R^1q_{C*}\mathbf{Z}_\ell(1) \to
R^{2n-1}f_*\mathbf{Z}_\ell(n)$ est nul --- ou, ce qui revient au même puisque ce
sont des faisceaux lisses sur une base connexe, si $Z\colon H^1(C_{\bar s},
\mathbf{Z}_\ell(1)) \to H^{2n-1}(Y_{\bar s}, \mathbf{Z}_\ell(n))$ est nul ---
alors $u(\mathrm{cl}(z)) = 0$.
\end{quote}

\pagerange{11}{12}
\subsection*{Décomposer un cycle algébriquement trivial (n\textsuperscript{o}~2, pages 11 et 12)}

Soit $i\colon Y \to X$ un morphisme de variétés propres et lisses sur un corps
algébriquement clos $k$, et
\[ i^*\colon H^{2n-1}(X, \mathbf{Q}_\ell(n)) \to H^{2n-1}(Y, \mathbf{Q}_\ell(n)). \]
On suppose :
\begin{itemize}
\item[a)] il existe $v\colon H^{2n-1}(Y) \to H^{2n-1}(X)$, induit par un cycle
algébrique (à coefficients rationnels) sur $Y \times X$, tel que $i^* v$ soit
l'identité sur l'image de $i^*$ ;
\item[b)] pour toute courbe propre, lisse et connexe $C$ sur $k$ et tout cycle
$Z$ de codimension $n$ sur $Y \times C$, l'image de
$Z\colon H^1(C, \mathbf{Q}_\ell(1)) \to H^{2n-1}(Y, \mathbf{Q}_\ell(n))$ est
contenue dans celle de $i^*$.
\end{itemize}
L'hypothèse a) est un cas particulier des conjectures standard ; elle est vraie
quand $i$ est l'inclusion d'une section hyperplane et que l'opérateur $\Lambda$
de Lefschetz est algébrique sur $X$, ce qu'on sait, dit le tapuscrit, quand $X$
se relève en caractéristique nulle avec sa polarisation. L'hypothèse b) dit,
dans le langage des motifs, que le plus grand morceau de niveau 1 du motif de
réalisation $\operatorname{Coker} i^*$ est nul ; en termes d'aujourd'hui, que
$\operatorname{Coker} i^*$ ne rencontre pas le cran de coniveau $n-1$ de la
cohomologie de $Y$.\footnote{« Niveau » est son mot, au sens qu'on dit
aujourd'hui « niveau de Hodge » : un morceau de $H^{2n-1}$ est de niveau 1 s'il
est engendré par des images de $H^1$ de courbes, c'est-à-dire s'il est dans le
cran $N^{n-1}$ de la filtration par le coniveau. « Plus grand » est ajouté à
l'encre.}

\emph{Énoncé.} Sous a) et b), tout cycle $z \in Z^n(Y)$ algébriquement équivalent
à zéro s'écrit, dans le groupe de Chow $\mathrm{CH}^n(Y) \otimes \mathbf{Q}$,
\[ z = i^*(z') + z'' , \]
où $z' \in Z^n(X)_{\mathbf{Q}}$ est algébriquement équivalent à zéro et
$z'' = Z''(t_1) - Z''(t_0)$ pour un $\mathbf{Q}$-cycle $Z''$ de codimension $n$
sur $Y \times C$ qui induit l'homomorphisme nul $H^1(C) \to H^{2n-1}(Y)$.

\emph{Démonstration.} On écrit $z = Z(t_1) - Z(t_0)$ avec $Z$ un cycle sur
$Y \times C$, $C$ une courbe lisse connexe.\footnote{Qu'une seule courbe suffise
à paramétrer un cycle algébriquement équivalent à zéro est classique ; le
tapuscrit le prend pour acquis.} Par b), $Z\colon H^1(C) \to H^{2n-1}(Y)$ tombe
dans l'image de $i^*$, donc par a) $Z = i^*(vZ)$ ; le composé
$vZ\colon H^1(C) \to H^{2n-1}(X)$ est induit par un $\mathbf{Q}$-cycle $Z'$ sur
$X \times C$, puisque $v$ l'est. On pose $Z'' = Z - (i \times
\mathrm{id}_C)^*Z'$, qui induit l'homomorphisme nul, puis $z' = Z'(t_1) -
Z'(t_0)$ et $z'' = Z''(t_1) - Z''(t_0)$.\footnote{Le tapuscrit écrit
« $(iv)Z = i(vZ)$ », $i$ pour $i^*$. Les égalités de cycles qu'il écrit sont des
égalités de classes : $(i \times \mathrm{id}_C)^*$ est le tiré en arrière du
calcul d'intersection sur les variétés lisses, défini modulo l'équivalence
rationnelle. D'où la formulation dans $\mathrm{CH}^n(Y) \otimes \mathbf{Q}$, et
les coefficients rationnels, que l'encre ajoute (« $\mathbf{Q}$- » au-dessus de
« cycle »).}

\pagerange{12}{13}
\subsection*{La version relative (n\textsuperscript{o}~3, pages 12 et 13)}

Soient maintenant $X$, $Y$ propres et lisses sur $S$ régulier et connexe,
$g\colon X \to S$, $f\colon Y \to S$, et $i\colon Y \to X$ un $S$-morphisme.
On suppose a) et b) vérifiés par $i_{\bar s}\colon Y_{\bar s} \to X_{\bar s}$,
où $\bar s$ est un point géométrique au-dessus du point générique de
$S$.\footnote{Le tapuscrit écrit « $i_{\bar s}\colon X_{\bar s} \to X_{\bar
s}$ ».} Soit $z \in Z^n(Y)$ dont la restriction à $Y_{\bar s}$ est
algébriquement, donc homologiquement, équivalente à zéro ; sa classe est
primitive et $u(z) \in H^1(S, R^{2n-1}f_*\mathbf{Z}_\ell(n))$ est définie.

\emph{Énoncé.} Quitte à remplacer $S$ par un ouvert non vide,
$u(z)$ appartient, modulo torsion, à l'image de
\[ i^*\colon H^1(S, R^{2n-1}g_*\mathbf{Z}_\ell(n)) \to H^1(S,
R^{2n-1}f_*\mathbf{Z}_\ell(n)) . \]
La restriction à un ouvert est, dit le tapuscrit, « probablement inutile ».

\emph{Démonstration.} On applique le n\textsuperscript{o}~2 à $i_{\bar s}$. La
courbe $C$, ses points $t_0$, $t_1$ et les cycles $Z'$, $Z''$ sont définis sur
une extension finie du corps des fonctions de $S$ ; un argument de trace ramène à
ce corps lui-même, au prix d'un multiple entier --- d'où le « modulo torsion » ---
et, sur un ouvert de $S$, $C$ devient une courbe relative propre et lisse, $Z''$
un cycle sur $Y \times_S C$, et la décomposition $z = i^*(z') + z''$ vaut sur
$S$. Alors $u(z) = i^*u(z') + u(z'')$, et $u(z'') = 0$ par le
n\textsuperscript{o}~1, puisque $Z''$ induit l'homomorphisme nul sur la fibre en
$\bar s$.

\pagerange{13}{15}
\subsection*{L'homomorphisme de Griffiths, et ce qu'il donne (n\textsuperscript{os}~4 et 5, pages 13 à 15)}

Soit $X$ projective et lisse de dimension $m$ sur $k$ algébriquement clos, et un
pinceau de Lefschetz ; l'éclatement de son axe donne $\widetilde{X}$ et
$\tilde f\colon \widetilde{X} \to \mathbf{P}^1_k$, dont les fibres sont les
sections du pinceau. Soit $S \subset \mathbf{P}^1_k$ l'ouvert où $\tilde f$ est
lisse, et $f\colon Y = \widetilde{X}|_S \to S$. La partie primitive $P^m(X)$ est
le noyau de $H^m(X) \to H^0(S, R^mf_*\mathbf{Q}_\ell)$, c'est-à-dire le noyau de
la restriction à une section hyperplane --- la partie primitive au sens de
Lefschetz. Le n\textsuperscript{o}~1 fournit l'\emph{homomorphisme de Griffiths}
\[ \gamma\colon P^m(X) \longrightarrow H^1(S, R^{m-1}f_*\mathbf{Q}_\ell). \]
L'immersion $i\colon Y \to X_S = X \times S$ au-dessus de $S$ donne
\[ i^*\colon H^1(S, H^{m-1}(X, \mathbf{Q}_\ell)) \to H^1(S,
R^{m-1}f_*\mathbf{Q}_\ell). \]

\begin{tikzcd}
Y \arrow[rr, "i"] \arrow[dr, "f"'] & & X_S \arrow[dl, "g"] \\
& S &
\end{tikzcd}

\emph{Énoncé.} $\gamma^{-1}(\operatorname{Im} i^*) = 0$. Le tapuscrit renvoie à
« des calculs explicites essentiellement triviaux, compte tenu de la structure
cohomologique connue d'une variété éclatée ». En voici la substance, qui est
nôtre : sous le théorème de Lefschetz difficile, le faisceau $R^{m-1}f_*\mathbf{Q}_\ell$
est somme directe de la partie constante $H^{m-1}(X)$ et de la partie
évanescente $\mathcal{E}$, l'image de
$i^*$ est le premier facteur de $H^1(S, -)$, et la composante de $\gamma$ dans
$H^1(S, \mathcal{E})$ est injective, parce que $P^m(X)$ s'identifie, par la suite
de Leray de $\tilde f$ et la théorie de Lefschetz, à $H^1(\mathbf{P}^1,
j_*\mathcal{E})$, qui s'injecte dans
$H^1(S, \mathcal{E})$.

\emph{Le théorème de Griffiths, conditionnel} (n\textsuperscript{o}~5). On
suppose $m = 2n$, et que l'inclusion $Y_{\bar s} \to X_{k(\bar s)}$ de la
section hyperplane générique satisfait a) et b). Soit $z \in Z^n(X)$ dont la
restriction à $Y_{\bar s}$ est algébriquement équivalente à zéro. Sa classe est
alors primitive ; par le n\textsuperscript{o}~3, $\gamma(\mathrm{cl}(z))$ tombe
modulo torsion dans l'image de $i^*$ ; par le n\textsuperscript{o}~4, elle est
nulle : \emph{$z$ est homologiquement équivalent à zéro}, à torsion près.

Soit $\operatorname{Grif}(Y_{\bar s})$ le groupe des cycles homologiquement
équivalents à zéro modulo ceux qui sont algébriquement équivalents à zéro --- le
nom est le sien, et c'est resté le nom. Restreindre à $Y_{\bar s}$ les cycles de
$X$ de classe primitive envoie ceux-ci dans $\operatorname{Grif}(Y_{\bar s})$,
et ce qui précède dit que le noyau de cette application, tensorisée par
$\mathbf{Q}$, est formé de cycles homologiquement triviaux sur $X$. Il en résulte
que \emph{$\operatorname{Grif}(Y_{\bar s}) \otimes \mathbf{Q}$ contient un
sous-espace qui se surjecte sur $P^{2n}_{\mathrm{alg}}(X)$}, l'espace des classes
primitives algébriques de $X$ --- espace qui, dit le tapuscrit, « a une nette
tendance à être $\neq 0$ ».\footnote{Le tapuscrit dit que $\operatorname{Grif}
\otimes \mathbf{Q}$ « contient comme sous-espace » $P^{2n}_{\mathrm{alg}}(X)$.
L'argument donne moins : un cycle homologiquement trivial de $X$ n'a aucune raison
de se restreindre en un cycle algébriquement trivial de $Y_{\bar s}$, de sorte
que la restriction n'est pas définie sur les classes de cohomologie. Ce qui est
établi, c'est un sous-quotient, et l'inégalité de dimensions qui suffit au
propos.} Dans l'exemple minimal de Griffiths --- $X$ quadrique de
$\mathbf{P}^5$, $z$ la différence de deux plans des deux familles --- c'est
exactement ce qu'il faut.

\pagerange{15}{16}
\subsection*{La conjecture du n\textsuperscript{o}~6 (pages 15 et 16)}

\emph{Conjecture.} Soit $X \subset \mathbf{P}^r_k$ comme au
n\textsuperscript{o}~5, avec $n \geqslant 2$. Il existe $d_0$ tel que, pour tout
pinceau de Lefschetz assez général d'hypersurfaces de degré $d \geqslant d_0$,
l'inclusion de la section générique $Y_{\bar s} \hookrightarrow X_{\bar s}$
satisfasse b).\footnote{« $\subset \mathbf{P}^r_k$ », « avec $n \geqslant 2$ »
et « de Lefschetz » sont tapés en interligne, « assez général » est à l'encre, et
l'énoncé est marqué en marge d'un double trait.} L'hypothèse a), elle, relève
de la conjecture d'algébricité de $\Lambda$.

On suppose le théorème de Lefschetz difficile vrai sur $X$ et $\Lambda$
algébrique.\footnote{« Lefschetz vache » est son mot pour le théorème de
Lefschetz difficile, conjectural en caractéristique $p$ quand ces lignes sont
écrites et démontré par Deligne (\emph{La conjecture de Weil II}, 1980).} Soit
$E^{2n-1} \subset H^{2n-1}(Y_{\bar s})$ l'orthogonal de $H^{2n-1}(X)$ : la
partie évanescente. La théorie des cycles évanescents donne qu'elle est
\emph{un module simple sous $\operatorname{Gal}(k(\bar s)/k(s))$ et sous chacun
de ses sous-groupes ouverts} ; le tapuscrit en voit la raison dans la parité :
en dimension impaire $2n-1$, les monodromies locales sont des transvections et
non des symétries, et aucune de leurs puissances ne s'annule.\footnote{L'argument
du tapuscrit est une indication. L'énoncé, sous sa forme « pour tout sous-groupe
ouvert », découle aujourd'hui de ce que l'adhérence de Zariski du groupe de
monodromie géométrique est le groupe symplectique tout entier, connexe, dès que
la monodromie n'est pas finie (Deligne, \emph{La conjecture de Weil I}, 1974,
§~5, via Kazhdan--Margulis). On signale le cas de monodromie finie comme une
réserve que le tapuscrit ne fait pas.} Comme la plus grande partie de niveau 1
est stable par Galois, la condition b) équivaut alors à : \emph{$E^{2n-1}$ n'est
pas tout entier de niveau 1}.

\emph{En caractéristique nulle,} si $E^{2n-1}$ était de niveau 1, ses
composantes de Hodge seraient nulles hors des types $(n, n-1)$ et $(n-1,
n)$\footnote{Le tapuscrit écrit ici $E^{n-1}$ ; c'est $E^{2n-1}$.}, donc
$H^{p,q}(X) \to H^{p,q}(Y)$ serait un isomorphisme pour les autres types, et en
particulier, comme $n \geqslant 2$, $H^{2n-1}(X, \mathcal{O}_X) \to
H^{2n-1}(Y, \mathcal{O}_Y)$ aussi. Or la dimension du second membre tend vers
l'infini avec $d$. Cela démontre la conjecture en caractéristique nulle, par la
méthode de Lefschetz reprise par Griffiths.

\emph{En caractéristique $p$,} on ne sait, dit-il, que des cas particuliers ;
Katz traite les intersections complètes et un pinceau assez général, ce qui suffit
à l'exemple minimal. Pour voir que $E^{2n-1}(Y_{\bar s})$ n'est pas de niveau 1, il
suffit de trouver un point $t$ du pinceau où $E^{2n-1}(Y_{\bar t})$ ne l'est pas ---
par le critère de bonne réduction de Néron--Ogg--Chafarevitch ; « la triste vérité
est que je ne sais pas montrer de façon générale, même en admettant les
conjectures standard, que le niveau d'un motif ne peut que diminuer par
spécialisation ». En se spécialisant sur un corps fini à $q$ éléments, le niveau 1
force les valeurs propres de Frobenius sur $E^{2n-1}(Y_{\bar t})$ à être divisibles
par $q^{n-1}$ (comme entiers algébriques) ; il suffit donc de trouver $t$ où l'une
d'elles est une unité $p$-adique. C'est ce que Katz obtient par la formule des
points fixes de Woods Hole et Verdier.\footnote{Le tapuscrit écrit
« Lefschetz-Woodshole-Verdier » : la formule des points fixes pour la cohomologie
cohérente, issue du colloque de Woods Hole (1964), sous la forme algébrique que
lui donne Verdier. « (et est peut-être équivalent à) », qui suit « implique »,
est ajouté à l'encre.}

\pagerange{16}{20}
\subsection*{Le théorème 7.1 et ses corollaires (n\textsuperscript{o}~7, pages 16 à 20)}

Le n\textsuperscript{o}~7 est écrit deux fois. La première rédaction tient en une
ligne, au bas de la page 16 : « Le procédé de Katz. A lui de jouer \ldots ». La
seconde, page 17, commence par « Je vais prouver la conjecture du
n\textsuperscript{o}~6 pour des pinceaux de Lefschetz assez généraux
d'hypersurfaces de grand degré », et l'encre entoure « vais prouver » et écrit
en tête du feuillet : « non, pas tout à fait ! ».\footnote{Le « non » est d'une
lecture incertaine ; « pas tout à fait » ne l'est pas.}

Dans tout ce numéro, $X \subset \mathbf{P}^r_k$ est projectif, équidimensionnel,
de Cohen--Macaulay, de dimension $m > 0$, sur un corps $k$ de caractéristique
$p > 0$.

\emph{Le cadre.} Soit $S_d$ l'espace projectif des formes de degré $d$, $k'$ son
corps des fonctions et $Y \subset X_{k'}$ la section de $X$ par l'hypersurface
générique de degré $d$.\footnote{Le tapuscrit écrit $k' = k(t_0, \ldots, t_r)$.
Les coefficients de la forme générique de degré $d$ sont indexés par les monômes
de degré $d$ en $r+1$ variables, et c'est en autant d'indéterminées que $k'$ est
engendré.} Soit $E_0(Y)$ le conoyau de $H^{m-1}(X_{k'}, \mathcal{O}) \to
H^{m-1}(Y, \mathcal{O}_Y)$, espace de dimension finie sur $k'$, muni du Frobenius
$p$-linéaire $F$ induit par l'élévation à la puissance $p$. Sa \emph{partie
semi-simple} $V$ est le plus grand sous-espace stable sur lequel $F$ est bijectif ;
l'ensemble $V_0$ des $x \in V \otimes \bar k'$ tels que $Fx = x$ est un
$\mathbf{F}_p$-espace de dimension $\dim_{k'}V$, sur lequel agit
$G = \operatorname{Gal}(\bar k'/k')$, d'image $G_0 \subset
\operatorname{Aut}(V_0)$. En langage d'aujourd'hui, $V_0$ est le
$\mathbf{F}_p$-système local étale qui correspond à la partie de pente nulle de
$(E_0, F)$, et $\dim V$ un rang de Hasse--Witt stable.

\emph{Théorème 7.1.} Il existe $d_0$ tel que pour tout $d \geqslant d_0$ et tout
$a \in \mathbf{F}_p$, il existe $u \in G_0$ et un sous-espace $V'_0 \subset V_0$
stable par $u$ avec $\operatorname{Tr}(u|V'_0) = a$.\footnote{C'est l'énoncé tel
que l'encre le corrige : « d'un sous-groupe $V'_0 \subset V_0$ stable par $G_0$ »
est inséré sous la ligne sans place marquée, et « $|V'_0$ » ajouté à la trace. La
stabilité que la démonstration fournit (lemme 7.5) est par $u$, non par $G_0$ ;
on écrit ce qui est démontré. Le texte tapé disait davantage : un $u \in G_0$ de
trace $a$ \emph{sur $V_0$ tout entier}, et ajoutait que $G_0$ n'est pas un
$p$-groupe --- parenthèse barrée à l'encre. Le tapuscrit écrit
$a \in \mathbf{F}_p^*$ ; la démonstration donne tout $a \in \mathbf{F}_p$.}

En marge de cet énoncé, à l'encre et de biais : « pas prouvé sous cette forme
juste ».\footnote{Lecture incertaine. Telle qu'elle se lit, et rapprochée de la
note du corollaire 7.2, elle signifie que l'énoncé tapé n'est pas prouvé ; on
verra que l'argument donne exactement la forme corrigée.}

\emph{Corollaire 7.2.} Supposons $k$ fini. Pour tout ouvert non vide $U$ de
$S_d$ et tout entier $a$, il existerait une extension finie $k'$ de $k$ et
$t \in U(k')$ tel que
\[ \operatorname{card} Y_t(k') \equiv a \pmod p , \]
où $Y_t$ est la section de $X_{k'}$ par l'hypersurface $t$. En marge, à l'encre :
« pas prouvé ; également 7.1 sous sa forme juste », et un « pas prouvé » au crayon
au bout de la démonstration.\footnote{La première marginale est d'une lecture
incertaine ; la seconde, dans un ovale au crayon, ne l'est pas. On écrit le
corollaire au conditionnel pour cette raison.}

La démonstration esquissée est la suivante. Pour $Y$ propre sur un corps fini, la
somme alternée des traces du Frobenius sur les $H^i(Y, \mathcal{O}_Y)$ est
congrue modulo $p$ au nombre de points rationnels --- la congruence que le
tapuscrit attribue à Deligne, par la formule de Lefschetz--Verdier.\footnote{Elle
est exposée, à notre connaissance, par Katz dans SGA~7~XXII, et étendue à tout
schéma propre par Fulton (1978).} Le choix de $d_0$ assure
$H^i(X, \mathcal{O}_X(-d)) = 0$ pour $i < m$ --- c'est là que sert l'hypothèse de
Cohen--Macaulay ---, donc $H^i(X, \mathcal{O}) \to H^i(Y_s, \mathcal{O})$ est un
isomorphisme pour $i \leqslant m-2$ et injectif pour $i = m-1$. Sur un ouvert
convenable $U$, la formation des $R^if_*\mathcal{O}$ commute au changement de
base, et $V$ s'étend en un sous-fibré $\underline{V}$ stable sur lequel $F$ est
bijectif ; la représentation $V_0$ provient alors du groupe fondamental $\pi$ de
$U$. On obtient, $f_t \in G_0$ désignant le Frobenius du point $t$ et $c_t$ la
somme alternée des traces de Frobenius sur les $H^i(X, \mathcal{O}_X)$,
$0 \leqslant i \leqslant m-1$,
\[ \operatorname{card} Y_t(k') \equiv c_t + (-1)^{m-1}\operatorname{Tr} f_t
\pmod p . \]
Quitte à agrandir $k$, $c_t = c$ ne dépend pas de $t$, et le théorème de densité
de Čebotarev donne un $t$ de Frobenius $f_t$ prescrit dans $G_0$.\footnote{Pour
les schémas de type fini sur un corps fini ; Serre, 1965. Le caron est ajouté à la
main.} Mais il faut pour cela un $f \in G_0$ de trace prescrite \emph{sur $V_0$
tout entier} : c'est la forme tapée de 7.1, que l'argument n'établit pas. Le
corollaire 7.2 reste donc non démontré, comme l'auteur le note.

\emph{Corollaire 7.3.} Supposons $m$ pair, $m \geqslant 4$, et $X$ lisse. Il
existerait un ouvert non vide $U'$ de la grassmannienne des droites de $S_d$ tel
que, pour toute droite $D \in U'(k')$ ($k'$ une extension quelconque de $k$)
définissant un pinceau de Lefschetz sur $X_{k'}$, la partie évanescente
$E(Y_t)$ de la section au point générique $t$ de $D$ ne soit pas de niveau 1, ni
ne contienne de sous-motif non nul de niveau 1, même après extension de
$k(t)$.\footnote{L'hypothèse « $m$ pair, $X$ lisse » est tapée, « $\geqslant 4$ »
et les ajouts sur $U'$ et $D$ sont à l'encre. La note encadrée de la marge définit
$U'$ comme un ouvert de droites rencontrant $\overline{U}$ telles que $\pi_1(D
\cap U) \to \pi_1(U)$ soit surjectif, et annonce « une variante de Bertini » ;
« ouvert » et « $U$ » y sont d'une lecture incertaine, et un mot y est illisible.}

L'argument : par le critère de Néron--Ogg--Chafarevitch, on se ramène à $k$ fini,
et il suffit de trouver un point fermé $s$ de $D$ où les valeurs propres de
Frobenius sur $E^{m-1}(Y_s, \mathbf{Q}_\ell)$ ne sont pas toutes divisibles par
$p$. Par la formule des traces, le théorème de Lefschetz faible et la dualité de
Poincaré, le nombre de points de $Y_s$ est la somme d'un terme venu de la
cohomologie de $X$ et de $(-1)^{m-1}\operatorname{Tr}(\mathrm{Frob} \mid
E^{m-1}(Y_s))$ ; les termes de degré $i \leqslant m-2$ et leurs duaux de degré
$2m-2-i$, une fois $k$ agrandi pour que les valeurs propres sur les
$H^i(X_{\bar k})$ soient congrues à $0$ ou $1$ modulo $p$, contribuent modulo
$p$ une constante $c$ indépendante de $s$ (les duaux sont divisibles par
$q$).\footnote{La formule de la page 20 est reprise à l'encre et coupée par le
bord du feuillet ; l'exposant du facteur $(1 + q^{\ldots})$ est illisible, et les
signes ont été biffés plusieurs fois. On donne la structure de la formule, non sa
lettre ; la justification par la divisibilité des termes duaux est nôtre.} Si
toutes les valeurs propres sur $E^{m-1}$ étaient divisibles par $p$, on aurait
$\operatorname{card} Y_s(k') \equiv c$ pour toute extension finie $k'$ ; il
suffit donc de trouver $s \in D \cap U$ et $k'$ avec $\operatorname{card}
Y_s(k') \not\equiv c$. C'est ce que donnerait 7.2, raffiné en imposant à $s$ de
rester sur une droite $D$ donnée : par Bertini, $\pi_1(D \cap U)$ se surjecte sur
$\pi_1(U)$, donc sur $G_0$, et Čebotarev s'applique sur $D \cap U$. Le
tapuscrit s'arrête là, au milieu de la page 20.

\pagerange{21}{26}
\subsection*{Remords, démonstration de 7.1, et le n\textsuperscript{o}~8 (pages 21 à 26)}

\emph{Remarques 7.4.} a) « Remords à 7.3 » : inutile de supposer $m$ pair et
$m \geqslant 4$ si l'on énonce que \emph{le niveau de $E^{m-1}(Y_t)$ est
$m-1$}, c'est-à-dire que $E^{m-1}(Y_t)$ n'a aucun morceau de coniveau
$\geqslant 1$. Point n'est besoin alors du critère de Néron--Ogg--Chafarevitch :
si un cycle sur $Z \times Y_{\bar t}$ définissait un épimorphisme
$H^{m-3}(Z)(-1) \to E(Y_{\bar t})$, il se spécialiserait aux points $s$ de $D$
voisins de $t$, et les valeurs propres de Frobenius sur $E(Y_{\bar s})$ seraient
divisibles par $\operatorname{card} k(s)$, ce qu'un choix de $D$ assez général
interdit. Si $m$ est impair, la représentation de monodromie ne reste plus
simple sur les sous-groupes ouverts, mais le plus grand sous-motif de niveau
$< m-1$ est stable par Galois, et la même conclusion vaut sur $k(t)$.\footnote{Le
tapuscrit tape « de niveau $n-1$ » en interligne, là où le sens demande « de
niveau $< n-1$ » --- la transcription le signale d'un « sic ? ». « épi » est tapé
au-dessus de « morphisme ».} b) Le langage des motifs n'a servi que de
commodité ; on peut l'éviter en définissant la filtration par le niveau sur la
cohomologie $\ell$-adique comme en a). « Ceci dit, l'énoncé 7.3 \ldots{} est
démontré sans aucune conjecture (pas même Lefschetz vache pour $X$ \ldots) » ---
\emph{une fois 7.2 acquis}, précise-t-on ici, puisque 7.3 en dépend.

\emph{Démonstration de 7.1.} On prend $d_0$ tel que $H^i(X,
\mathcal{O}_X(-d)) = 0$ pour $d \geqslant d_0$ et $i < m$, et pour $U \subset
S_d$ l'ouvert des $s$ où $\dim Y_s = m-1$, qui est aussi le plus grand ouvert où
$\underline{Y} \to S_d$ est plat. Par EGA~III~7, les $R^if_*\mathcal{O}_Y$ sont
plats sur $U$ et commutent au changement de base. « On fera attention » : on
devra utiliser des points $s$ où $Y_s$ est réductible, et ce n'est qu'une fois
7.1 acquis qu'on pourrait descendre dans un petit ouvert de $S_d$, comme en 7.2.
On se ramène à $k$ fini à $q = p^f$ éléments, et, quitte à l'agrandir, à ce que
les valeurs propres des Frobenius sur les $H^i(X, \mathcal{O}_X)$ soient $0$ ou
$1$ ; soit $c$ la somme alternée des rangs semi-simples des $H^i(X,
\mathcal{O}_X)$, $0 \leqslant i \leqslant m-1$. La congruence de Deligne et les
isomorphismes précédents donnent, pour tout $s \in U(k')$,
\[ \operatorname{card} Y_s(k') \equiv c + (-1)^{m-1}\operatorname{Tr} f_s
\pmod p , \]
où $f_s$ est le Frobenius relatif à $k'$ sur la partie semi-simple de $E_0(Y_s)$.
\footnote{Le tapuscrit note $q = p^m$ ; on écrit $p^f$, $m$ étant la dimension.
L'encre corrige « c'est la puissance $m$-ème du frobenius » en « ce sont les
puissances \ldots{} des frobenius ».} Le théorème découle de deux lemmes.

\emph{Lemme 7.5.} Pour tout $s \in U$, il existe $f \in G_0$ et un sous-espace
$V'_0 \subset V_0$ stable par $f$ tels que le polynôme caractéristique de $f$
sur $V'_0$ soit celui de $f_s$ ; en particulier $\operatorname{Tr}(f|V'_0) =
\operatorname{Tr} f_s$.\footnote{L'ajout « et $V'_0 \subset V_0$ stable par
$f$ » est à l'encre ; « stable par » est d'une lecture incertaine.}

\emph{Lemme 7.6.} Pour tout entier $N \geqslant 0$, il existe $d'_0$ tel que pour
tout $d \geqslant d'_0$ il existe une extension finie $k'$ de $k$ et $s \in
U_d(k')$ avec $\operatorname{card} Y_s(k') = N$. Cela vaut dès que $X$ est un
sous-schéma fermé de $\mathbf{P}^r_k$ de dimension $m \geqslant 1$.

\emph{Démonstration de 7.6.} On peut supposer $X$ réduit. Soit $e$ la somme des
degrés des composantes irréductibles de dimension $m$, et $d_1$ tel que
$N \leqslant e\, d_1^{\,m}$.\footnote{Le tapuscrit écrit $r$ pour cette somme, et
« $N \geqslant r d_1^{\,n}$ ». L'inégalité doit aller dans l'autre sens : on veut
découper $N$ points parmi les $e\,d_1^{\,m}$ points de $Z$. Le tapuscrit le
confirme lui-même quelques lignes plus bas, où le cas $d_1 = 1$ est celui de
« $N \leqslant r$ ».} Quitte à agrandir $k$, les composantes $X_i$ sont
géométriquement irréductibles, chacune a un point rationnel $x_i$, et $m$ formes
$\phi_0, \ldots, \phi_{m-1}$ de degré $d_1$ découpent sur $X$ un schéma fini
étale $Z$ de degré $e\,d_1^{\,m}$, à points tous rationnels, évitant les $x_i$.
Quitte à agrandir encore $k$ et $d_1$, une forme $\phi_m$ de degré $d_1$ découpe
sur $Z$ exactement $N$ points.\footnote{Ce pas est affirmé, non argumenté.} Pour
$d_2 \geqslant m+1$, il existe sur un corps fini une forme $\underline{\phi}$ en
$m+1$ variables, de degré $d_2$, sans zéro non trivial --- une forme normique,
restreinte au besoin à un sous-espace. Alors $\phi =
\underline{\phi}(\phi_0, \ldots, \phi_m)$, de degré $d_1d_2$, a pour zéros
rationnels sur $X$ exactement les $N$ points de $Z \cap H_{\phi_m}$, et ne
s'annule en aucun $x_i$, donc définit un point de $U$. Si $N \leqslant e$, on
prend $d_1 = 1$, et $d'_0 = m+1$ convient. Sinon, on multiplie $\phi$ par une forme
$\psi$ de degré $d_3 \geqslant m+1$ sans zéro rationnel sur $X$ (le cas $N = 0$),
ce qui ne change pas l'ensemble des zéros rationnels, et l'on atteint tous les
degrés $d = d_1d_2 + d_3$ ; avec $d_2 = m+1$, on peut prendre
\[ d'_0 = (d_1 + 1)(m+1). \]
\footnote{Le tapuscrit écrit $(H_{\phi\psi} \cap X)(k) = Z(k)$ ; c'est
$(H_{\phi_m} \cap Z)(k)$, de cardinal $N$.}

\emph{Comment 7.1 en résulte.} Pour $a \in \mathbf{F}_p$, on choisit $N \in [0,
p-1]$ congru à $c + (-1)^{m-1}a$, on prend $d \geqslant d_0$ et $\geqslant$ les
$d'_0(N)$ ; 7.6 donne $s$ avec $\operatorname{Tr} f_s = a$, et 7.5 un $f \in G_0$
et un $V'_0$ stable par $f$ avec $\operatorname{Tr}(f|V'_0) = a$. C'est la forme
corrigée de 7.1, et rien de plus : le passage par un sous-espace est inévitable,
parce que les points $s$ fournis par 7.6 sont des sections réductibles, dont la
partie semi-simple peut être plus petite que la partie générique.

\emph{Remarques 7.7.} a) On obtient ainsi une estimation de $d_0$ en fonction de
$p$, de $e$ et du plus petit entier qui assure l'annulation des $H^i(X,
\mathcal{O}(-d))$. Si l'on se contente d'au moins deux traces distinctes, on peut
prendre $d'_0 = m+1$ ; cela suffit, dans la forme corrigée, à établir que
$V \neq 0$, c'est-à-dire que le Frobenius n'est pas nilpotent sur
$E_0(Y)$.\footnote{Le tapuscrit en déduit aussi que $G_0$ n'est pas un
$p$-groupe : cela vaut pour la forme tapée de 7.1 (les éléments d'un $p$-groupe
de $\operatorname{GL}(V_0)$ ont tous la trace $\dim V_0$), non pour la forme
corrigée, où des sous-espaces différents donnent des traces différentes même si
$G_0$ est trivial. Il en déduit enfin que ce choix de $d_0$ suffit à une forme
affaiblie de 7.2, donc à 7.3 : cela suppose de nouveau la forme tapée.} b) On
peut améliorer $d'_0$ en découpant $N = N' + N''$ et en prenant des formes
produits de degrés premiers entre eux.\footnote{Le tapuscrit écrit « $N''
\geqslant r d_2'^{\,n}$ » ; le sens demande $d''_1$, et l'inégalité dans le même
sens qu'en 7.6, corrigée comme plus haut.} c) Les formes auxiliaires sans zéro
sont des formes normiques, géométriquement produits de formes linéaires : c'est
pourquoi on ne peut travailler que dans le gros ouvert $U_d$.

\emph{Lemme 7.8 (la raison de 7.5).} Soit $S$ un schéma normal de
caractéristique $p$, $\underline{V}$ un module localement libre de type fini muni
d'un endomorphisme $p$-linéaire, $t$ un point de $S$ et $s$ une spécialisation de
$t$. Soient $V_1$ et $V_0$ les $\mathbf{F}_p$-représentations de $G_t$ et $G_s$
définies par les parties semi-simples des fibres en $t$ et en $s$. Il existe un
sous-groupe de décomposition $D \subset G_t$ au-dessus de $s$ et un sous-espace
$V'_0 \subset V_1$ stable par $D$, tels que l'action de $D$ sur $V'_0$ se factorise
par $D \to G_s$ et soit isomorphe à $V_0$.\footnote{Le tapuscrit dit « un
sous-groupe \emph{ouvert} $G'_t$ de $G_t$ » ; un groupe de décomposition est
fermé et n'est pas ouvert en général. Ce qui suffit à 7.5, c'est qu'il contienne
un relèvement du Frobenius de $s$.} La démonstration, dit-il, est « par l'âne qui
trotte », en se ramenant à $S$ intègre et hensélien, de point fermé $s$ et de point
générique $t$. Elle tient en ceci, qui est nôtre : l'équation $Fx = x$ définit un
revêtement étale, de sorte que ses solutions dans la fibre spéciale se relèvent de
façon unique au-dessus du hensélisé, et s'injectent dans la fibre générique. C'est
la version en caractéristique $p$ de ce que le rang de Hasse--Witt stable (le
$p$-rang, pour une variété abélienne) ne peut que baisser par spécialisation.

Le tapuscrit s'achève, page 26, sur un titre à l'encre sans texte :
« 8. Relations heuristiques entre cohomologie cristalline et la torsion
cristalline, et la cohomologie cohérente. (J'ai donné notes à
Berthelot) ».\footnote{« torsion » et « Berthelot » sont d'une lecture
incertaine.}

\pagerange{27}{33}
\section*{Biextensions de complexes $[M \to G]$ (pages 27 à 33)}

Notes à l'encre, écrites pour trois d'entre elles au dos de feuillets du
tapuscrit de l'exposé VIII. Elles posent la situation que Deligne publiera en 1974
sous le nom de \emph{1-motifs} : une variété abélienne
munie d'un réseau et d'un tore, et sa duale.\footnote{Deligne, \emph{Théorie de
Hodge III}, Publ.\ IHES 44 (1974), §~10. Ces feuillets ne sont pas datés ; on ne
dit rien de l'ordre entre les deux, et le nom de « 1-motif » n'est pas sur la
page.}

\emph{La situation} (pages 27 et 28). Soient $B$, $B'$ deux schémas abéliens
duaux, $\varepsilon$ la biextension de Poincaré de $B \times B'$ par
$\mathbb{G}_m$, $M$, $M'$ deux groupes « réseaux » (localement libres de type
fini), et $i\colon M \to B$, $i'\colon M' \to B'$. Comme
$\underline{\operatorname{Hom}}(B, \mathbb{G}_m) = 0$ et
$\underline{\operatorname{Ext}}^1(M', \mathbb{G}_m) = 0$, on a
\[ \operatorname{Hom}(M', \underline{\operatorname{Ext}}^1(B, \mathbb{G}_m))
\simeq \operatorname{Biext}(B, M'; \mathbb{G}_m) \simeq
\operatorname{Ext}^1(B, D(M')) , \]
où $D(M') = \underline{\operatorname{Hom}}(M', \mathbb{G}_m)$ est le tore dual de
$M'$ ; comme $\underline{\operatorname{Ext}}^1(B, \mathbb{G}_m) = B'$,
\emph{se donner $i'$ revient à se donner une extension}
\[ 0 \to D(M') \to G \to B \to 0 , \]
et symétriquement $i$ donne $0 \to D(M) \to G' \to B' \to 0$. On pose $T =
D(M')$, $T' = D(M)$. Les \emph{relèvements} $M \to G$ de $i$ correspondent alors
aux trivialisations de la biextension $\varepsilon_{i,i'}$ tirée en arrière sur
$M \times M'$, et de même les relèvements $M' \to G'$ de $i'$ ; et les deux
trivialisations de $\varepsilon_{M,M'}$ qu'on déduit d'un côté et de l'autre
coïncident.\footnote{L'« Exemple » de la page 28, d'où viennent ces équivalences,
est d'une lecture incertaine par endroits (« si », « vrai », « ainsi »,
« Donc »), et un passage biffé y est à peu près illisible. Les deux conditions
d'annulation sont écrites sous les flèches, entourées.}

\emph{Le complexe et sa biextension} (page 30). Posons $A = [M \to G]$ et $A' =
[M' \to G']$, complexes en degrés $-1$ et $0$. Lorsque $M \to G$ et $M' \to G'$
sont injectifs, on a une biextension canonique de $(G/M, G'/M')$ ; en général, on
dispose d'une classe canonique
\[ \xi \in \operatorname{Ext}^1(A \overset{L}{\otimes} A', \mathbb{G}_m), \]
qui est la classe d'une biextension de $(A, A')$ par $\mathbb{G}_m$ (SGA~7~VII).
Le groupe de $n$-torsion du complexe, ${}_nA = \underline{\operatorname{Ext}}^0
(\mathbf{Z}/n, A)$, s'insère dans
\[ 0 \to {}_nG \to {}_nA \to M_n \to 0 \]
dès que ${}_nM = 0$ et que $G$ est $n$-divisible\footnote{La page écrit
« ${}_nG = G$ », puis « i.e.\ $nB = B$, $nT = T$ » : c'est la divisibilité qui
est voulue. Pour un groupe extension d'un schéma abélien par un tore, elle est
d'ailleurs automatique en topologie plate. La parenthèse « ce dernier résulte de
${}_nM = 0$ » ne se comprend pas telle quelle.}, et à la limite sur $n =
\ell^\nu$
\[ 0 \to T_\ell(G) \to T_\ell(A) \to M \otimes \mathbf{Z}_\ell \to 0 , \]
et de même pour $A'$.\footnote{Le premier terme de la seconde suite se lit
$T_\ell(A')$ ; on lit $T_\ell(G')$ par symétrie, comme la transcription.} C'est
la filtration par le poids du module de Tate d'un 1-motif : $T_\ell(T) \subset
T_\ell(G) \subset T_\ell(A)$.

\emph{L'accouplement et l'orthogonalité} (page 32). La classe $\xi$ induit
\[ \varphi^\xi_n\colon {}_nA \times {}_nA' \to \mu_n , \qquad
T_\ell(A) \times T_\ell(A') \to \mathbf{Z}_\ell(1) . \]

\emph{Proposition.} Dans cet accouplement, ${}_nT = D(M'_n)$ est orthogonal à
${}_nG'$, et ${}_nG$ est orthogonal à ${}_nT' = D(M_n)$. Les accouplements induits
--- entre ${}_nA/{}_nG = M_n$ et ${}_nT'$, entre ${}_nA'/{}_nG' = M'_n$ et
${}_nT$, et entre ${}_nG/{}_nT = {}_nB$ et ${}_nG'/{}_nT' = {}_nB'$ --- sont
« ceux qu'on devine » : les dualités de Cartier évidentes et l'accouplement de
Weil.\footnote{Sous $D(M'_n)$, la page écrit « $= {}_nT'$ », et sous $D(M_n)$
« $= {}_nT$ », ce qui intervertit les identifications des pages 27 et 28
($T = D(M')$, $T' = D(M)$) ; on suit ces dernières. Les deux parenthèses « (le
reste si ${}_nM' = 0$) » sont d'une lecture incertaine.}

\emph{Corollaire.} Si ${}_nM = {}_nM' = 0$ et que l'accouplement de Weil entre
${}_nB$ et ${}_nB'$ est parfait, l'accouplement entre ${}_nA$ et ${}_nA'$ est
parfait --- le gradué l'est.\footnote{Pour des schémas abéliens duaux,
l'accouplement de Weil sur les ${}_n$ est toujours parfait (dualité de Cartier
entre ${}_nB$ et ${}_nB'$) ; l'hypothèse est donc automatique. L'exemple que la
page en donne est d'une lecture incertaine.}

\emph{Les feuillets tapés} (pages 29, 31, 33). Au dos des pages 28, 30 et 32, et
reproduits tête-bêche, trois feuillets d'un état dactylographié de l'exposé VIII
de SGA~7, §~2, « Accouplements définis par une biextension » (de 2.1 à la fin de
2.3) : un texte publié, que la transcription ne reproduit pas. On y voit ses
corrections à l'encre --- des renvois rectifiés, des flèches nommées dans les
diagrammes, et, à la place de « qu'on laisse au lecteur le soin d'établir »,
« cf.~2.3 ci-dessous pour une démonstration ».

\pagerange{34}{45}
\section*{Picard et dualité (pages 34 à 45)}

Une couverture porte ce titre (page 34). Les feuillets qui suivent tournent autour
d'une seule formule --- les torseurs sur $X$ sous un groupe commutatif $G$ se
lisent comme des homomorphismes du dual de $G$ vers le foncteur de Picard de $X$ ---
et de ses variantes.\footnote{Rien sur les feuillets 41 à 45 ne dit qu'ils
appartiennent tous à la chemise « Picard et dualité » ; on les y range parce
qu'ils en continuent les formules.} La topologie est la topologie plate (fppf)
partout où interviennent $\mu_p$ et $\alpha_p$.

\pagerange{35}{35}
\subsection*{La cohomologie plate d'une courbe (page 35)}

Une table, pour $X$ une courbe propre, lisse et connexe sur $k$ algébriquement
clos de caractéristique $p$, de jacobienne $P$.\footnote{La page ne dit pas ces
hypothèses ; ce sont celles où la table est vraie.}
\[
\begin{array}{c|ccc}
 & H^0 & H^1 & H^2 \\
\hline
\mu_\ell & \mu_\ell(k) & {}_\ell P & \mathbf{Z}/\ell \\
\mu_p & 0 & {}_pP(k) & \mathbf{Z}/p \\
\alpha_p & 0 & \operatorname{Ker}\bigl(F \mid H^1(X, \mathcal{O}_X)\bigr) &
\operatorname{Coker}\bigl(F \mid H^1(X, \mathcal{O}_X)\bigr)
\end{array}
\]
avec, pour les coefficients constants, $H^i(X, \mathbf{Z}/\ell) = H^i(X,
\mu_\ell) \otimes \operatorname{Hom}(\mu_\ell(k), \mathbf{Z}/\ell)$, et
$H^0(X, \mathbf{Z}/p) = \mathbf{Z}/p$, $H^1(X, \mathbf{Z}/p) = H^1(X,
\mathcal{O}_X)^{F=1}$, $H^2(X, \mathbf{Z}/p) = 0$. Les lignes $\mu$ sont la
théorie de Kummer ; les lignes $p$ sont la dualité de Serre--Cartier, que la marge
nomme ; l'accouplement ${}_\ell P \times {}_\ell P \to \mu_\ell$ est l'autodualité
de la jacobienne.\footnote{Il note d'un $z$ bouclé ce qu'on écrit $\mu_\ell$,
$\mu_p$, et $U_\ell$ le groupe $\mu_\ell(k)$ ; il note $i_p$ ce qu'on écrit
$\alpha_p$, noyau du Frobenius de $\mathbb{G}_a$. La page ne nomme pas ces
faisceaux : leur identification est nôtre, et repose sur ce que la table leur
attribue exactement leurs valeurs. Elle écrit « $H^2(X, \mathbf{Z}/p) \simeq
H^2(X, \mathbf{Z}/p) = 0$ » tel quel, et « et cohomologie discrète », d'une
lecture incertaine, en marge.}

\pagerange{36}{36}
\subsection*{Une adjonction, et un programme (page 36)}

Pour $f\colon X \to S$ et $B$ un faisceau sur $S$,
\[ \operatorname{Ext}_X(f^*B, \mathbb{G}_m) \simeq \operatorname{Ext}_S(B,
Rf_*\mathbb{G}_m) , \]
d'où une suite spectrale d'aboutissement $\operatorname{Ext}_X$ et de terme
$H^p(X, \underline{\operatorname{Ext}}^q_X(f^*B, \mathbb{G}_m))$.\footnote{La page
écrit $f^!$ ; l'adjoint à droite de $f^*$ est $Rf_*$, et c'est lui que la formule
demande ($f^!$ est l'adjoint à droite de $Rf_!$). Même lecture plus bas pour ses
$f_!$ et $f_{\cdot}$.} Suit, en accolade, une liste de sujets, presque tous
d'une lecture incertaine : suites exactes de cohomologie et critères de
représentabilité ; Frobenius et Verschiebung ; dualité de Cartier ; un théorème
de finitude ; propriétés formelles, puis différentielles, des groupes et des
fibrés ; et le cas particulier des fibrés linéaires, symplectiques, orthogonaux.
C'est un plan, et le dossier n'en développe que le début.

\pagerange{38}{40}
\subsection*{Torseurs sous une variété abélienne (pages 38 à 40)}

Soit $X$ propre et lisse sur $k$ algébriquement clos, avec $H^0(X,
\mathcal{O}_X) = k$, et $A$ une variété abélienne sur $k$, de duale
$A'$.\footnote{La page dit « $X$ propre et simple sur $k$ », annonce des
« généralisations à une base quelconque », et ne dit pas que $k$ est
algébriquement clos ; c'est ce qu'il faut pour que $A(k)$ soit divisible et que les
torseurs sous un groupe fini venus de la base soient triviaux.} On sait que
$H^1(X, A)$ est de torsion, et la suite de Kummer $0 \to {}_nA \to A \to A \to 0$
donne
\[ 0 \to A(X)_n \to H^1(X, {}_nA) \to {}_nH^1(X, A) \to 0 . \]
On a $A(X)_n = \operatorname{Hom}(\operatorname{Alb}_X, A)_n$, et
$H^1(X, G) \simeq \operatorname{Hom}(D(G), \underline{\operatorname{Pic}}_X)$
pour un groupe commutatif fini $G$, de dual de Cartier $D(G)$ ; comme
$\operatorname{Hom}(\operatorname{Alb}_X, A) \simeq \operatorname{Hom}(A',
\underline{\operatorname{Pic}}^0_{X,\mathrm{red}})$ (la duale de l'albanaise) et $D({}_nA) \simeq {}_nA'$,
\[ (*) \qquad 0 \to \operatorname{Hom}(A', \underline{\operatorname{Pic}}^0_X)_n \to
\operatorname{Hom}({}_nA', \underline{\operatorname{Pic}}_X) \to {}_nH^1(X, A) \to 0 . \]
D'autre part, la suite $0 \to {}_nA' \to A' \to A' \to 0$, à laquelle on
applique $\operatorname{Hom}(-, \underline{\operatorname{Pic}}_X)$, donne
\[ (**) \qquad 0 \to \operatorname{Hom}(A', \underline{\operatorname{Pic}}_X)_n \to
\operatorname{Hom}({}_nA', \underline{\operatorname{Pic}}_X) \to
{}_n\operatorname{Ext}^1(A', \underline{\operatorname{Pic}}_X) \to 0 . \]
Les premiers termes coïncident ($A'$ est connexe), et à la limite inductive sur $n$
\[ 0 \to \operatorname{Hom}(A', \underline{\operatorname{Pic}}_X) \otimes
\mathbf{Q}/\mathbf{Z} \to \operatorname{Hom}(T_{\cdot}(A'),
\underline{\operatorname{Pic}}_X) \to H^1(X, A) \to 0 , \]
où $\operatorname{Hom}(T_\cdot(A'), -) = \varinjlim_n \operatorname{Hom}({}_nA',
-)$, et de même avec $\operatorname{Ext}^1$. D'où
\[ \boxed{H^1(X, A) \simeq \operatorname{Ext}^1(A',
\underline{\operatorname{Pic}}_X)_{\mathrm{tors}}} \]
à comparer à la formule de dualité $H^1(X, G) \simeq \operatorname{Hom}(D(G),
\underline{\operatorname{Pic}}_X)$ pour $G$ fini.\footnote{Deux corrections. La
page écrit $\otimes_{\mathbf{Z}} \mathbf{Q}$ ; la limite inductive des $N/nN$
pour les flèches de multiplication est $N \otimes \mathbf{Q}/\mathbf{Z}$. Et la
limite des ${}_n\operatorname{Ext}^1$ est le sous-groupe de torsion de
$\operatorname{Ext}^1$, non $\operatorname{Ext}^1$ tout entier comme l'écrit
l'encadré : les deux coïncident si $\operatorname{Ext}^1(A',
\underline{\operatorname{Pic}}_X)$ est de torsion, ce que la page ne discute pas.
Il faut enfin que les flèches de $(*)$ et $(**)$ se correspondent, ce qu'elle
admet. L'exposant $0$ du dernier terme de $(**)$ est d'une lecture incertaine.}
En marge : « Cet isomorphisme reste-t-il valable pour $X$ singulier ??? Il est
possible qu'on ait besoin de $H^i(X, \mathbb{G}_m) = 0$ pour $i \geqslant 2$ ».
Une remarque introduit les flèches canoniques
\[ \operatorname{Ext}^1(\underline{\operatorname{Alb}}_X, A) \to H^1(X, A), \qquad
\operatorname{Ext}^1(\underline{\operatorname{Alb}}_X, A) \simeq
\operatorname{Ext}^1(A', \underline{\operatorname{Alb}}_X{}') \simeq
\operatorname{Ext}^1(A', \underline{\operatorname{Pic}}^{\tau}_X) , \]
dont il affirme qu'on vérifie que ce sont des isomorphismes sur une partie que la
page ne laisse pas lire.\footnote{L'exposant $\tau$ est d'une lecture incertaine ;
les trois dernières lignes de la page 39 ne sont lues qu'en partie.}

Page 40, la même question en relatif, pour $X \to S$ et un $S$-groupe
$\gamma$, écrite avec le faisceau $\underline{\mathcal{H}}^1(X/S; -) =
R^1f_*$ :
\[ R^1f_*\underline{\operatorname{Hom}}_{S\text{-gr}}(\gamma, G) \simeq
\underline{\operatorname{Hom}}_{S\text{-gr}}(\gamma, R^1f_*G), \qquad
R^1f_*\underline{\operatorname{Ext}}^1_{S\text{-gr}}(\gamma, G) \simeq
\underline{\operatorname{Ext}}^1_{S\text{-gr}}(\gamma, R^1f_*G) , \]
la première donnée comme isomorphisme de foncteurs « quand tout est affine sur
$S$ (?) », la seconde suivie de « Sous quelles conditions ??? ». Le feuillet ne
tranche pas, et l'on ne tranche pas pour lui.

\pagerange{41}{43}
\subsection*{Un morphisme birationnel, et l'adjonction pour $\mathbb{G}_m$ (pages 41 à 43)}

La page 41, barrée de deux diagonales, considère un morphisme propre et
birationnel $f\colon X \to Y$ de schémas intègres, et l'ouvert $U \subset Y$ où
$f^{-1}$ est défini ; si $Y$ est régulier en codimension 1, le complémentaire de
$U$ est de codimension $\geqslant 2$. Une classe $\xi \in \operatorname{Pic}(X)$
donne par la section $g\colon U \to f^{-1}(U)$ une classe de
$\operatorname{Pic}(U)$, donc un cycle divisoriel sur $Y$ --- et la page s'arrête
là.

La page 42 écrit l'adjonction pour $f\colon X \to Y$ :
\[ \boxed{\underline{\operatorname{Ext}}^i_X(f^*G, H) \simeq
\underline{\operatorname{Ext}}^i_Y(G, Rf_*H)} , \]
puis la spécialise à $H = \mathbb{G}_m$, où $Rf_*\mathbb{G}_m$ est le complexe
qu'il note $\underline{\operatorname{Pic}}^\bullet_{X/Y}$, de faisceaux de
cohomologie $f_*\mathbb{G}_m$, $\underline{\operatorname{Pic}}_{X/Y} =
R^1f_*\mathbb{G}_m$, $R^2f_*\mathbb{G}_m$, \ldots{} --- « les supérieurs ».
\footnote{La page écrit $f_!$ et $f_{\cdot}$ ; voir la page 36. Les exposants
$(1)$, $(2)$ de $\underline{\operatorname{Pic}}$ sont d'une lecture incertaine ;
on les lit comme les $R^1$, $R^2$.} Pour $G$ fini et plat, de dual $D(G)$, avec
$f_*\mathcal{O}_X = \mathcal{O}_Y$ et $Y$ le spectre d'un corps algébriquement
clos, il en sort une suite spectrale
\[ E_2^{p,q} = \operatorname{Ext}^p(G, R^qf_*\mathbb{G}_m) \Rightarrow
H^{p+q}(X, D(G)) , \]
dont il écrit les bas degrés : $H^0(X, D(G)) = D(G)$, $H^1(X, D(G)) =
\operatorname{Hom}(G, \underline{\operatorname{Pic}}_X)$ --- c'est l'encadré
$H^1(X, G) \simeq \operatorname{Hom}(D(G), \underline{\operatorname{Pic}}_X)$, aux
rôles de $G$ et $D(G)$ près ---, et, pour $i = 2$,
\[ H^2(X, D(G)) \neq \operatorname{Ext}^1(G, \underline{\operatorname{Pic}}_X)
\oplus \operatorname{Hom}(G, R^2f_*\mathbb{G}_m) \quad ! \]
Le signe $\neq$ et le point d'exclamation sont les siens : $H^2$ n'est pas une
somme, mais porte une filtration dont le gradué est $\operatorname{Ext}^1(G,
\underline{\operatorname{Pic}}_X)$ et le noyau de la différentielle
$\operatorname{Hom}(G, R^2f_*\mathbb{G}_m) \to \operatorname{Ext}^2(G,
\underline{\operatorname{Pic}}_X)$.\footnote{Cette lecture du « $\neq$ ! » est
nôtre. La dernière ligne de la page, sur $H^{i-1}(X, G')$, est lue en partie, et
ses exposants sont illisibles.}

La page 43 commence une définition et l'abandonne : « \emph{Foncteurs adjoints.}
Définition 1. Soient $X, Y \in (\mathrm{Cat})$, $f\colon X \to Y$ et $g\colon Y
\to X$, on dit que $(f, g)$ est un couple de foncteurs adjoints, ou que $g$ est
adjoint de $f$, ou que $f$ est coadjoint de $g$, s'il existe un \ldots
».\footnote{La phrase s'arrête avant de dire de quel côté est l'adjonction ; on ne
peut donc pas traduire « adjoint » et « coadjoint » en « adjoint à droite » et « à
gauche ».} Suivent la curryfication
\[ \operatorname{Hom}(X, \underline{\operatorname{Hom}}_{S\text{-gr}}(\gamma,
G)) \simeq \{X \times \gamma \to G, \text{ additifs en } \gamma\} \simeq
\operatorname{Hom}_{S\text{-gr}}(\gamma, \underline{\operatorname{Hom}}_S(X, G)) ,
\]
et, pour un $S$-groupe commutatif $\gamma$ et $X$ muni d'une section $a$, la suite
exacte scindée par l'évaluation en $a$
\[ 0 \to \gamma \to \underline{\operatorname{Hom}}_S(X, \gamma) \to
\operatorname{Hom}((X, a), (\gamma, e)) \to 0 , \]
avec $\underline{\operatorname{Hom}}_S(X, \gamma) = \gamma$ lorsque
$f_*\mathcal{O}_X = \mathcal{O}_S$ et $\gamma$ est affine.\footnote{La page écrit
« $\doteq \gamma$ », sans dire sous quelles hypothèses ; on donne celles où c'est
vrai. Une égalité entre $H^1(X, \underline{\operatorname{Hom}}(\gamma, G))$ et
$H^1(X, \underline{\operatorname{Ext}}^1(\gamma, G))$, écrite en tête, est
abandonnée aussitôt, et on ne la reprend pas.}

\pagerange{44}{45}
\subsection*{Le groupe fondamental d'une variété abélienne (pages 44 et 45)}

Soit $A$ une variété abélienne sur $k$ algébriquement clos de caractéristique
$p$, $A^*$ sa duale. Le fait qui porte tout est qu'\emph{un revêtement étale
d'une variété abélienne est une variété abélienne} (Serre--Lang) : le groupe
fondamental est abélien, limite des noyaux des isogénies.

\emph{La partie première à $p$} (page 45). $\pi_1^{(p')}(A) \simeq
\varprojlim_{(n,p)=1} {}_nA$, et pour $\ell \neq p$
\[ \pi_1^{\ell}(A) \simeq T_\ell(A) = \operatorname{Hom}(\mathbf{Q}_\ell/
\mathbf{Z}_\ell, A). \]
Par Kummer, $\operatorname{Hom}(\pi_1^\ell(A), \mu_{\ell^\infty})$ est le
sous-groupe de $\operatorname{Pic}(A) = H^1(A, \mathcal{O}_A^*)$ formé des
éléments d'ordre une puissance de $\ell$. Si le groupe de Néron--Severi de $A$ est
sans torsion, ce sous-groupe est $A^*[\ell^\infty]$, et l'on obtient
l'accouplement de Weil
\[ \boxed{T_\ell(A) \times T_\ell(A^*) \longrightarrow T_\ell(\mu) =
\mathbf{Z}_\ell(1)} , \]
parfait, et ses niveaux finis ${}_{\ell^k}A \times {}_{\ell^k}A^* \to
\mu_{\ell^k}$.\footnote{« Si on admet que $A$ \ldots{} pas de torsion » : le mot
illisible est sans doute le sujet de « pas de torsion », et la lecture de ce
dernier est incertaine. L'hypothèse dont l'argument a besoin est que
$\mathrm{NS}(A)$ soit sans torsion, ce qui est vrai pour une variété abélienne
($\underline{\operatorname{Pic}}^\tau_A = \underline{\operatorname{Pic}}^0_A$).
La page note $U_\ell$ le groupe des racines de l'unité d'ordre une puissance de
$\ell$.}

\emph{La partie $p$} (page 44). Par la théorie d'Artin--Schreier--Witt, on a
\[ H^1(A, \mathbf{Z}/p^n) \simeq H^1(A, W_n\mathcal{O}_A)^{F=1} , \]
donc
\[ \operatorname{Hom}(\pi_1^p(A), \mathbf{Z}_p) = H^1(A, \mathbf{Z}_p) \simeq
H^1(A, W\mathcal{O}_A)^{F=1} , \]
et, comme $\pi_1^p(A) \simeq T_p(A)$, un accouplement parfait de
$\mathbf{Z}_p$-modules libres de type fini
\[ \boxed{T_p(A) \times H^1(A, W\mathcal{O}_A)^{F=1} \longrightarrow
\mathbf{Z}_p} , \]
de rang le $p$-rang de $A$ ; après extension des scalaires,
$\operatorname{Hom}_{\mathrm{cont}}(T_p(A), W(k)) \simeq H^1(A,
W\mathcal{O}_A)_s$, la partie de pente nulle.\footnote{Il écrit
$\mathbb{W}$ pour $W\mathcal{O}_A$, l'exposant $(p)$ pour les points fixes de $F$,
et l'indice $s$ pour la partie semi-simple. Le $\Lambda$ d'un encadré est d'une
lecture incertaine ; on lit $W(k)$, comme dans le dernier.}

\pagerange{46}{58}
\section*{Biextensions de faisceaux de groupes (pages 46 à 58)}

Une couverture au crayon porte ce titre (page 46). Toutes les biextensions sont
de faisceaux abéliens sur un topos, par un faisceau abélien $G$.

\pagerange{47}{50}
\subsection*{Reconstruire une biextension à partir d'une restriction (pages 47 à 50)}

Soit $0 \to P' \to P \to P'' \to 0$ exacte, et $Q$, avec a) $Q$ $n$-divisible et
b) ${}_nP \to P''$ épimorphique.\footnote{La page écrit « a) ${}_nQ = Q$ » ; le
sens, et le corollaire 1 qui écrit « $pQ = Q$ », demandent la divisibilité.} Une
biextension $E$ de $P \times Q$ par $G$ définit a) sa restriction $E'$ à
$P' \times Q$, d'où un accouplement $E'(n)\colon {}_nP' \times {}_nQ \to G$ ; b)
sa restriction à ${}_nP \times Q$, c'est-à-dire un accouplement $\varphi\colon
{}_nP \times {}_nQ \to G$ ; et $\varphi$ prolonge $E'(n)$. Ces accouplements sont
à valeurs dans ${}_nG$, par bilinéarité.

« Je dis que le foncteur $E \mapsto (E', \varphi)$ est une équivalence de
catégories. » La pleine fidélité revient à
$\underline{\operatorname{Hom}}(P'' \otimes Q, G) = 0$, ce que a) et b)
assurent puisqu'ils donnent $P'' \otimes Q = 0$ ; le reste de l'argument est
barré de diagonales et n'est pas repris.\footnote{Le passage barré invoque une
condition c) qui n'est pas sur la page.}

\emph{Corollaire 1.} Soit $\mathbb{P}$ un ensemble de nombres premiers ; on
suppose $\varinjlim_n {}_nP \to P''$ surjectif, $n$ parcourant les entiers à
facteurs dans $\mathbb{P}$, et $pQ = Q$ pour $p \in \mathbb{P}$. Alors $E
\mapsto (E', \varphi)$ est une équivalence entre les biextensions de $(P, Q)$
par $G$ et les couples formés d'une biextension $E'$ de $(P', Q)$ et d'un
accouplement
\[ \varphi\colon T_{\mathbb{P}}(P) \times T_{\mathbb{P}}(Q) \to
T_{\mathbb{P}}(G) \]
prolongeant celui, $\varphi_{\mathbb{P}, E'}$, que définit $E'$. Un « NB »
encadré, lu en partie, précise qu'il faut entendre les $T_{\mathbb{P}}$ comme
des systèmes projectifs et garder le système des ${}_nP \times {}_nQ \to {}_nG$
avec leurs compatibilités, plutôt que leur limite.

\emph{Corollaire 2.} Soient $0 \to P' \to P \to P'' \to 0$ et $0 \to Q' \to Q
\to Q'' \to 0$ exactes ; on suppose a) $pP = P$ et $pQ' = Q'$ pour $p \in
\mathbb{P}$ (ou la condition symétrique), b) $\varinjlim {}_nP \to P''$ et
$\varinjlim {}_nQ \to Q''$ épimorphiques. Alors $E \mapsto (E', \varphi)$, où
$E'$ est la restriction à $P' \times Q'$ et $\varphi\colon T_{\mathbb{P}}(P)
\times T_{\mathbb{P}}(Q) \to T_{\mathbb{P}}(G)$ prolonge $\varphi_{E',
\mathbb{P}}$, est une équivalence de catégories. On applique le corollaire 1 deux
fois : à $0 \to P' \to P \to P'' \to 0$ et $Q'$, puis à $P$ et $0 \to Q' \to Q \to
Q'' \to 0$.

Les deux corollaires reposent sur l'équivalence de la page 47, dont la
démonstration est barrée ; le dossier ne les établit pas.

Les pages 48 et 50 sont deux feuillets, numérotés 15 et 13, d'un tapuscrit
français qui n'est pas de lui et n'appartient pas à SGA~7 --- préfaisceaux et
faisceaux de modules sur un site, comparaison entre cohomologie de Čech et
cohomologie des faisceaux ---, avec des corrections à l'encre d'une main non
établie.

\pagerange{51}{53}
\subsection*{Biextensions sans hypothèse de rigidité (pages 51 et 53)}

Soient $A$, $G$ des faisceaux abéliens au-dessus de $S$, $f\colon A \to S$
le morphisme structural, $e_A$ la section unité, $G_A = f^*G$.

a) Si $\underline{\operatorname{Hom}}_{\mathrm{gr}}(A, G) = 0$, la suite
spectrale locale-globale donne
\[ \operatorname{Ext}^1(A, G) \simeq H^0(S,
\underline{\operatorname{Ext}}^1(A, G)) : \]
les faisceaux d'extensions déterminent les extensions.

b) Si $\underline{\operatorname{Hom}}_{\mathrm{pct}}(A, G) = 0$ --- les morphismes
de faisceaux d'ensembles qui respectent la section unité, de sorte que $G
\xrightarrow{\sim} f_*G_A$ ---, la suite de Leray, scindée par $e_A$, donne
\[ H^1(A, G_A) \simeq H^1(S, G) \times H^0(S, R^1f_*G_A) , \]
et $H^0(S, R^1f_*G_A)$ classifie les $G_A$-torseurs rigidifiés le long de
$e_A$. Le foncteur qui envoie une extension de $A$ par $G$ sur le torseur
rigidifié sous-jacent est pleinement fidèle, d'où une injection
$\operatorname{Ext}^1(A, G) \hookrightarrow H^1(A, G_A)$, d'image essentielle
les torseurs $L$ tels que
\[ \Lambda(L) = m^*L \otimes \mathrm{pr}_1^*L^{-1} \otimes \mathrm{pr}_2^*L^{-1}
\]
soit trivial, $m$ désignant l'addition. On a donc une suite exacte de faisceaux,
et une flèche vers le Néron--Severi :

\begin{tikzcd}[column sep=small]
0 \arrow[r] & \underline{\operatorname{Ext}}^1(A, G) \arrow[r, "i"] & R^1f_*G_A \arrow[r, "j"] \arrow[d, "k"'] & \underline{\operatorname{Tors}}_{\mathrm{birig,sym}}(A, A; G) \\
 & & \underline{\mathrm{NS}}(A) \arrow[ur, "j'"'] &
\end{tikzcd}

où $j(L) = \Lambda(L)$ et où $j'$ dit que les bitorseurs symétriques
birigidifiés de cette forme sont des biextensions.\footnote{La page dessine $k$
en oblique et met $j'$ sur une seconde ligne, sous un signe d'égalité. Le
« pct » est glosé en marge : « Hom (pas nécessairement multiplicatifs) qui
respectent la section unité » ; « pas nécessairement » est d'une lecture
incertaine. Le premier membre de la formule de $H^1$ est écrit $A_G$ ; on lit
$G_A$, comme au second.} Pour $G = \mathbb{G}_m$ et $A$ un schéma abélien, c'est
la situation classique : $\underline{\operatorname{Ext}}^1(A, \mathbb{G}_m) = A'$,
$R^1f_*\mathbb{G}_m = \underline{\operatorname{Pic}}_{A/S}$, et $j'$ envoie une
classe de Néron--Severi $L$ sur le fibré de Mumford $\Lambda(L)$, c'est-à-dire sur
l'homomorphisme symétrique $\varphi_L\colon A \to A'$.

c) Soit $B$ un autre faisceau abélien. Si
$\underline{\operatorname{Hom}}_{\mathrm{gr}}(A, G) = 0$, la catégorie
$\underline{\operatorname{BIEXT}}(A, B; G)$ est discrète, et
\[ \operatorname{Biext}(A, B; G) \simeq \operatorname{Hom}_{\mathrm{gr}}(B,
\underline{\operatorname{Ext}}^1(A, G)) . \]
Si l'on suppose seulement $\underline{\operatorname{Hom}}_{\mathrm{pct}}(A, G) =
0$, on a des injections
\[ \operatorname{Biext}(A, B; G) \hookrightarrow
\operatorname{Hom}_{\mathrm{gr}}(B, R^1f_*G_A) \hookrightarrow H^0(B,
R^1f_{B*}G_{A \times_S B}) , \]
le dernier terme classifiant les $G_{A \times_S B}$-torseurs rigidifiés le long
de $e_{A \times_S B}$. Les torseurs qui sont des biextensions sont ceux qui sont
rigidifiables, et pour lesquels l'application $B \to R^1f_*G_A$ qu'ils
définissent est additive et se factorise par $\underline{\operatorname{Ext}}^1(A,
G)$ ; la marge rappelle que rigidifiabilité et factorisation ensemble disent que
l'application symétrique $A \to R^1q_*G_B$ est additive. On obtient des foncteurs
pleinement fidèles de catégories discrètes
\[ \underline{\operatorname{BIEXT}}(A, B; G) \to
\underline{\operatorname{TORS\,BIRIG}}_{e_A, e_B}(A, B; G) \to
\underline{\operatorname{TORS\,RIG}}_{e_{A \times_S B}}(G_{A \times_S B}) , \]
qui sont, dit la page, des équivalences « dans deux cas importants » qu'elle ne
nomme pas.\footnote{« pl.\ fid. », « équivalence » et « vérifiable » sont d'une
lecture incertaine dans cette dernière phrase.}

\pagerange{55}{57}
\subsection*{La diagonale, et le groupe de Néron--Severi (pages 55 et 57)}

On définit une flèche en sens inverse de $j$ en tirant en arrière par la
diagonale $\Delta\colon A \to A \times_S A$ :
\[ \Delta^*\colon \underline{\operatorname{Tors}}_{\mathrm{birig,sym}}(A, A; G) \to
R^1f_*G_A . \]
Comme $m \circ \Delta = 2\,\mathrm{id}_A$ et $\mathrm{pr}_i \circ \Delta =
\mathrm{id}_A$,
\[ \Delta^* j(L) \simeq (2\,\mathrm{id}_A)^*(L) \otimes L^{-2} , \]
et modulo $A' = \underline{\operatorname{Ext}}^1(A, G)$, c'est-à-dire dans le
Néron--Severi, cela vaut $4L - 2L = 2L$. Cela résulte du

\emph{Lemme.} Pour tout entier $n$, $(n\,\mathrm{id}_A)^*L \equiv n^2L$ dans
$\underline{\mathrm{NS}}(A)$. Plus généralement, $u \mapsto u^*$, de
$\operatorname{Hom}(A, B)$ dans $\operatorname{Hom}(\underline{\mathrm{NS}}(B),
\underline{\mathrm{NS}}(A))$, est quadratique :
\[ \varphi(u+v) + \varphi(u-v) = 2\bigl(\varphi(u) + \varphi(v)\bigr) .
\]
\footnote{La page porte, d'une lecture incertaine, $\varphi(u+v) - \varphi(u-v)$.
Avec le signe $-$, $v = 0$ donnerait $0 = 2\varphi(u)$ ; c'est l'identité du
parallélogramme qui est vraie, et qui donne, pour $u = v$, $\varphi(2u) =
4\varphi(u)$.} En effet $u \mapsto (u \times u)^*$ est quadratique sur les
biextensions, par biadditivité, et la flèche des biextensions vers
$\underline{\mathrm{NS}}$ est linéaire.

Page 57 : le composé $k \circ \Delta^* j\colon \underline{\mathrm{NS}}_{A/S} \to
\underline{\mathrm{NS}}_{A/S}$ est la multiplication par $2$. Par conséquent
\emph{l'image réciproque par $2\,\mathrm{id}$ de l'extension canonique}
\[ 0 \to A' \to \underline{P}_{A/S} \to \underline{\mathrm{NS}}_{A/S} \to 0 \]
\emph{est canoniquement scindée} par $\Delta^* j$ ; autrement dit la classe de
l'extension est tuée par $2$, et l'extension est canoniquement scindée après
inversion de $2$.\footnote{La page dit que l'extension « splitte canoniquement »,
avec deux mots illisibles avant « de l'extension » ; « par » et « splitte » sont
d'une lecture incertaine, et elle écrit « par $P_{A/S}$ » là où le noyau de
$k\colon \underline{P}_{A/S} \to \underline{\mathrm{NS}}_{A/S}$ est $A'$. Un
relèvement $s$ avec $ks = 2$ ne scinde pas l'extension elle-même ; on écrit ce
qu'il scinde.} Il commence ensuite à dire que la restriction de $j \circ \Delta^*$
aux bitorseurs symétriques est la multiplication par $2$ --- ce qui est vrai, une
biextension symétrique $E$ vérifiant $\Lambda(\Delta^*E) \simeq E \otimes
\sigma^*E \simeq E^2$ ---, et la page s'arrête.

Les pages 52, 54, 56 et 58 sont des feuillets d'une rédaction anglaise de la
théorie géométrique des invariants de Mumford. Le seul qui porte quelque chose de
lui, la page 56, a en marge des remarques au crayon, en anglais, présumées de sa
main, sur les hypothèses d'une proposition : « Seems to hold for any equivalence
relation \ldots{} Is \emph{flat} enough ?? », « Doesn't seem enough, $G$
\emph{smooth}/$k$ ! », « ambiguous ».

\pagerange{59}{70}
\section*{Le théorème des cycles invariants pour le $H^1$ (pages 59 à 70)}

Une couverture, de sa main : « "Th.\ du cycle invariant" : Cas du $H^1$ » (page
59). Un tapuscrit corrigé à l'encre et au crayon (pages 60 à 63) énonce le
théorème et trois corollaires ; des notes manuscrites (pages 64 à 70) le
démontrent dans le cas d'un entier premier à $p$, d'abord sous l'hypothèse d'un
zéro-cycle de degré 1, puis sans elle.

Le problème est celui du \emph{cycle invariant} : pour une famille sur un disque,
les classes de la fibre générale fixes par la monodromie viennent-elles de
l'espace total ? Le théorème des cycles invariants le dit en caractéristique
nulle, et ce dossier en donne, pour le $H^1$ et au-dessus d'un trait, la forme
arithmétique exacte, avec le défaut qu'il mesure.\footnote{En $\ell$-adique, le
théorème des cycles invariants général est démontré par Deligne
(\emph{La conjecture de Weil II}, 1980) sous des hypothèses de géométrie
(familles sur une courbe sur un corps fini) qui ne sont pas celles d'ici.}

\pagerange{60}{63}
\subsection*{L'énoncé (pages 60 à 63)}

Soit $S$ un trait strictement local (le spectre d'un anneau de valuation discrète
strictement hensélien $V$), de point fermé $s$ de corps $k$ de caractéristique $p$,
de point générique $\eta$ de corps $K$, et $I = \operatorname{Gal}(K^s/K)$, qui
est ici le groupe d'inertie. Soit $f\colon X \to S$ projectif, $X$ régulier, avec
$\mathcal{O}_S \xrightarrow{\sim} f_*\mathcal{O}_X$, de sorte que $X$ est plat sur
$S$ et que la fibre spéciale $X_0$ est un diviseur. On note $(D_i)_{i \in J}$ ses
composantes irréductibles réduites et
\[ X_0 = \sum_{i \in J} d_iD_i, \qquad d = \operatorname{pgcd}(d_i) . \]
Soit $\Delta \simeq \mathbf{Z}^J$ le groupe des diviseurs de support dans $X_0$, et
\[ \Gamma = \Delta / \mathbf{Z}X_0 \hookrightarrow \operatorname{Pic}(X)
\]
son image dans $\operatorname{Pic}(X)$\footnote{La page dit « la catégorie des
diviseurs » ; c'est le groupe. Dans $\Gamma \simeq {}/\mathbf{Z}.X_0$, le
numérateur est laissé en blanc par la machine. Que $\Delta/\mathbf{Z}X_0$
s'injecte dans $\operatorname{Pic}(X)$ vient de ce qu'une fonction dont le
diviseur est porté par $X_0$ est inversible sur $X_\eta$, donc dans $K^*$ puisque
$H^0(X_\eta, \mathcal{O}) = K$.}, et
\[ w\colon \Gamma \to \operatorname{Pic}(X) \to \operatorname{Pic}(X_0) \to
\mathrm{NS}(X_0) = \operatorname{Im}\bigl(\operatorname{Pic}(X_0) \to
\underline{\mathrm{NS}}_{X_0/k}(k)\bigr) . \]
Soit enfin $d^t$ l'indice de $X_\eta$, le pgcd des degrés de ses zéro-cycles.

\emph{Théorème.}
\begin{itemize}
\item[a)] $d^t = \operatorname{pgcd}(\mu_id_i)$, où $\mu_i$ est la multiplicité
radicielle de $D_i$ sur $k$. En particulier $d$ divise $d^t$ et $d^t = p^hd$ pour
un $h \geqslant 0$ ; si $k$ est parfait, $h = 0$ et $d^t = d$.\footnote{L'encre
remplace la fin de la phrase tapée par « on a $d^t = p^hd$, où $p^h =
\operatorname{pgcd}(\mu_i)$ ». La première moitié est juste, la seconde non :
$p^h$ est seulement divisible par $\operatorname{pgcd}(\mu_i)$. Avec deux
composantes, $d_1 = 1$, $\mu_1 = p$ et $d_2 = p$, $\mu_2 = 1$, on a $d = 1$,
$\operatorname{pgcd}(\mu_i) = 1$ et $d^t = p$. La marginale « donc $h = 0$ i.e.\
$d = d_t$ si $k$ est parfait » est juste. La formule a) elle-même est, dit le
tapuscrit, « essentiellement une reformulation d'un résultat de M.~Raynaud », dont
la référence est laissée en blanc ; le renvoi à EGA~IV l'est aussi.}
\item[b)] $w$ est un homomorphisme de groupes abéliens de type fini \emph{dont le
noyau est contenu dans le sous-groupe de torsion de $\Gamma$}. Ce dernier est
canoniquement $\mathbf{Z}/d$, engendré par l'image de $\frac1d X_0 = \sum
\frac{d_i}{d}D_i$. Donc $R = \operatorname{Ker} w$ est cyclique d'ordre $d'$
divisant $d$, et a fortiori $d^t$.
\item[c)] Supposons $n$ premier à $p$, ou $d^t$ premier à $p$, ou $k$ parfait
(donc algébriquement clos). On a une suite exacte canonique, fonctorielle en $X$,
\[ (6) \qquad 0 \to {}_n\Gamma \to H^1(X, \mu_n) \xrightarrow{u_n}
H^1(X_{\bar\eta}, \mu_n)^I \to \operatorname{Ker} w_n \to 0 , \]
où $w_n = w \otimes \mathbf{Z}/n\colon \Gamma_n \to \mathrm{NS}(X_0)_n$ et
${}_n\Gamma \simeq \mathbf{Z}/(n, d)$ ; elles forment un système projectif quand
$n$ varie.
\item[d)] Pour tout nombre premier $\ell$, on a une suite exacte canonique,
fonctorielle en $X$,
\[ (7) \qquad 0 \to H^1(X, \mathbf{Z}_\ell(1)) \xrightarrow{u(\ell)}
H^1(X_{\bar\eta}, \mathbf{Z}_\ell(1))^I \to R \otimes \mathbf{Z}_\ell \to 0, \]
avec $R \otimes \mathbf{Z}_\ell = 0$ si $\ell \nmid d'$ --- a fortiori si $\ell
\nmid d^t$ --- et $\mathbf{Z}/\ell$ si $\ell \mid d'$.
\end{itemize}
\footnote{Pour $n$ ou $\ell$ divisible par $p$, il s'agit de cohomologie plate.
« fonctorielle en $X$ » est ajouté au crayon ; une parenthèse sur le passage à la
limite, sous les conditions de c), est encadrée et barrée au crayon. Le $\ne$ tapé
de « $\ell \ne d'$ » est surchargé en $\nmid$.}

\emph{Corollaire 9.} $u(\ell)$ est injectif, de conoyau fini, nul pour presque
tout $\ell$ ; de même avec $\mathbf{Z}_\ell$ au lieu de $\mathbf{Z}_\ell(1)$ si
$\ell \neq p$ --- les racines de l'unité d'ordre premier à $p$ étant dans $K$,
l'inertie agit trivialement sur $\mathbf{Z}_\ell(1)$.

\emph{Corollaire 10.} Pour $n$ quelconque, on a
\[ 0 \to {}_n\Gamma \to H^1(X, \mu_n) \xrightarrow{u_n} H^1(X_{\bar\eta},
\mu_n)^I \to Q(n) \to 0, \qquad 0 \to \operatorname{Ker} w_n \to Q(n) \to
\Theta(n) \to 0 , \]
avec
\[ \Theta(n) \hookrightarrow {}_n\operatorname{Ker}\bigl(\mathrm{Br}(K) \to
\mathrm{Br}(X_\eta)\bigr) \hookrightarrow {}_{(n, d^t)}\mathrm{Br}(K) ,
\]
\footnote{Le tapuscrit écrit $\mathrm{Br}(X_{\bar\eta})$. Le noyau de
$\mathrm{Br}(K) \to \mathrm{Br}(X_{\bar\eta})$ est $\mathrm{Br}(K)$ tout entier,
puisque cette flèche se factorise par $\mathrm{Br}(\bar K) = 0$ ; c'est le noyau
vers $\mathrm{Br}(X_\eta)$ qui contient l'obstruction et qui est tué par l'indice
$d^t$.} de sorte que $\Theta(n) = 0$ si $n$ ou $d^t$ est premier à $p$ (le groupe
de Brauer de $K$ est de $p$-torsion), ou si $k$ est parfait ($\mathrm{Br}(K) =
0$). En tous cas, si $p^r$ est la plus grande puissance de $p$ divisant $d^t$,
$p^r\Theta(n) = 0$ pour tout $n$. Comme les flèches de transition sur les
$\Theta(\ell^\nu)$ sont des multiplications par $\ell$, $\varprojlim_\nu
\Theta(\ell^\nu) = 0$, et (7) se déduit de (11) par passage à la
limite.\footnote{Pour $\ell = p$, ce passage à la limite demande une condition
de Mittag-Leffler sur les $H^1(X, \mu_{p^\nu})$ plats, que le tapuscrit ne
vérifie pas.}

\emph{Corollaire 15.} Écrivons $d^t = p^ra$ avec $(a, p) = 1$. Il existe un
diagramme, unique à isomorphisme près,

\begin{tikzcd}
X \arrow[d, "f"'] & X' \arrow[l, "h"'] \arrow[d, "f'"] \\
S & S' \arrow[l, "g"']
\end{tikzcd}

où $S' \to S$ est le revêtement kummérien connexe (modérément ramifié) de groupe
$\mu_a$, $h$ un revêtement étale, le carré cartésien au-dessus de
$\eta$\footnote{La page porte une lettre tracée à la main en forme de $\eta$ devant
« $\to S$ », sur un blanc de la machine ; on lit le changement de base $S' \to
S$, restreint au point générique.} ; $f'$ satisfait aux mêmes hypothèses que $f$,
et la partie première à $p$ de $d(X'/S')$ est $1$.\footnote{La page dit
« $d(X'/S')$ est maintenant égal à $p^r$ ». Le revêtement de degré $a$ divise les
multiplicités par $a$, qui est la partie première à $p$ de $d$ ; $d(X'/S')$ est la
partie $p$ de $d$, qui vaut $p^r$ quand $d^t = d$ (par exemple $k$ parfait), et
qui peut être plus petite sinon (voir a)). Seule compte, pour la suite, qu'elle
soit une puissance de $p$. Que $h$ soit étale vient de ce que $a$ divise tous les
$d_i$ et est premier à $p$.} Si $I' \subset I$ est le groupe d'inertie de $S'$,
pour tout $n$ premier à $p$,
\[ u'_n\colon H^1(X', \mu_n) \to H^1(X'_{\bar\eta}, \mu_n)^{I'} \]
est injectif, de conoyau ${}_n\delta/{}_n\delta'$, où $\delta$ est le
sous-groupe de torsion de $\operatorname{Coker} w$ et $\delta'$ celui de
$\mathrm{NS}(X_0)$, pris pour $X'$.\footnote{La page renvoie à « $w$ défini dans
(5) », c'est-à-dire pour $X$ ; c'est pour $X'$ que le calcul des pages 64 à 66
donne ce conoyau.} Par suite, pour tout $\ell \neq p$,
\[ u'(\ell)\colon H^1(X', \mathbf{Z}_\ell) \xrightarrow{\sim} H^1(X'_{\bar\eta},
\mathbf{Z}_\ell)^{I'} \]
\emph{est un isomorphisme}. Le tapuscrit s'arrête sur cet énoncé, au milieu de la
page 63 ; aucune démonstration ne suit dans le dossier.

\pagerange{64}{68}
\subsection*{La démonstration, pour $n$ premier à $p$ et un zéro-cycle de degré 1 (pages 64 à 68)}

La page 64 reprend l'énoncé à la main, avec $X$ propre et l'hypothèse que $X_\eta$
a un zéro-cycle de degré 1. Soit $\delta$ le sous-groupe de torsion du conoyau de
$\Delta \to \underline{\mathrm{NS}}_{X_0/k}(k)$, $\delta'$ celui de
$\mathrm{NS}(X_0)$, $\delta'' = \delta/\delta'$. Pour $n$ premier à $p$ on a
\[ 0 \to H^1(X, \mu_n) \xrightarrow{u_n} H^1(X_{\bar\eta}, \mu_n)^I \to
{}_n\delta/{}_n\delta' \to 0 , \]
et, à la limite, pour $\ell \neq p$,
\[ H^1(X, \mathbf{Z}_\ell) \xrightarrow{\sim} H^1(X_{\bar\eta},
\mathbf{Z}_\ell)^I \]
--- « Théorème des cycles invariants ».\footnote{La suite de la page 64 se
prolonge, d'une lecture douteuse, par « $\to {}_n\delta'' \to {}_n\delta' \to
\delta_n \to \delta''_n \to 0$ ». Le terme qui s'insère après $u_n$ est
${}_n\delta/{}_n\delta'$, sous-groupe de ${}_n\delta''$, comme le montre la suite
exacte du serpent que la page 68 écrit pour $0 \to \delta' \to \delta \to \delta''
\to 0$ ; on le donne. Le terme ${}_n\Gamma$ que la page écrit en tête est nul ici,
$\Gamma$ étant sans torsion (voir plus bas).}

\emph{Démonstration.} $X$ est intègre et plat sur $S$, donc $H^1(X, \mu_n)
\hookrightarrow H^1(X_\eta, \mu_n)$. Par Kummer, et parce que $V^*/V^{*n} = 0$ et
$K^*/K^{*n} = \mathbf{Z}/n$ ($V$ strictement local, $n$ premier à $p$),
\[ H^1(X, \mu_n) \simeq {}_n\operatorname{Pic}(X), \quad H^1(X_\eta,
\mu_n)/(\mathbf{Z}/n) \simeq {}_n\operatorname{Pic}(X_\eta), \quad
H^1(X_{\bar\eta}, \mu_n) \simeq {}_n\operatorname{Pic}(X_{\bar\eta}) . \]
Le zéro-cycle de degré 1 donne $\operatorname{Pic}(X_\eta) \simeq
\underline{\operatorname{Pic}}_{X_\eta/K}(K) = P(K) = P(K^s)^I$. Ainsi $u_n$
s'identifie à ${}_n\operatorname{Pic}(X) \to {}_n\operatorname{Pic}(X_\eta)$.
Comme $X$ est régulier --- il suffirait que ses anneaux locaux soient factoriels
---, on a
\[ 0 \to \Gamma \xrightarrow{i} \operatorname{Pic}(X) \xrightarrow{v}
\operatorname{Pic}(X_\eta) \to 0 , \]
d'où, par le lemme du serpent, $\operatorname{Ker} u_n = {}_n\Gamma$ et
$\operatorname{Coker} u_n \simeq \operatorname{Ker}(i_n\colon \Gamma_n \to
\operatorname{Pic}(X)_n)$.

\emph{Lemme} (marge de la page 66). $\operatorname{Pic}(X)_n \hookrightarrow
\operatorname{Pic}(X_0)_n$, du moins si $S$ est complet, ce qui suffit : les
obstructions à relever d'un voisinage infinitésimal au suivant sont dans des
$k$-vectoriels, où $n$ est inversible.\footnote{La page renvoie à « PCA, exposé
sur Picard » et dit le lemme « facile » ; la justification est résumée.} Donc
$\operatorname{Ker} i_n = \operatorname{Ker} j_n$, $j\colon \Gamma \to
\operatorname{Pic}(X_0)$. Puis $\operatorname{Pic}(X_0) =
\underline{\operatorname{Pic}}_{X_0/k}(k)$, le groupe
$\operatorname{Pic}^0(X_0)$ est $n$-divisible, et $\operatorname{Pic}(X_0)_n =
\mathrm{NS}(X_0)_n$. On obtient $\operatorname{Coker} u_n \simeq
\operatorname{Ker} w_n$.

Il dit alors que \emph{$\Gamma$ est sans torsion et $w$ injectif} (le lemme
marqué (D), page 68, « ici $X$ régulier, semble-t-il »). Que $\Gamma$ soit sans
torsion : l'adhérence schématique dans $X$ d'un point fermé de $X_\eta$ est finie
et plate sur $S$, de degré égal à sa multiplicité d'intersection avec $X_0 =
\sum d_iD_i$, donc multiple de $d$ ; ainsi $d \mid d^t$, et $d^t = 1$ force $d =
1$.\footnote{La page écrit « $\Gamma = \mathbf{Z}/d\mathbf{Z}$ » pour le
sous-groupe de torsion de $\Gamma$.} Avec $0 \to \Gamma \to \mathrm{NS}(X_0) \to
\operatorname{Coker} w \to 0$, le serpent donne $\operatorname{Ker} w_n \simeq
\operatorname{Coker}({}_n\mathrm{NS}(X_0) \to {}_n\operatorname{Coker} w) =
{}_n\delta/{}_n\delta'$, l'application $\delta' \to \delta$ étant injective parce
que $\Gamma$ est sans torsion. À la limite sur les $\ell^\nu$, ces groupes finis
s'annulent et l'on obtient l'isomorphisme.

Que $w$ soit injectif modulo torsion --- l'assertion b) ---, la page 68 le renvoie
à Raynaud : « Or c'est fait par Raynaud dans sa note : spécialisation des
foncteurs de Picard, th.~3 ». Une tentative barrée voulait se ramener à la
dimension relative 1, où c'est la négativité de la matrice d'intersection des
$D_i$.\footnote{M.~Raynaud, « Spécialisation du foncteur de Picard », publié in
extenso aux Publ.\ IHES 38 (1970) ; « note » et « th.~3 » sont la page. En
dimension relative 1, le noyau de la forme d'intersection sur $\mathbf{Q}^J$ est
la droite engendrée par $X_0$ (lemme de Zariski) : c'est b).}

\pagerange{70}{70}
\subsection*{Variante : sans zéro-cycle de degré 1 (page 70)}

« Laissons tomber l'hypothèse » $d^t = 1$.\footnote{La page écrit « l'hypothèse $d'$ ($=$ pgcd des deg.\ des $0$-cycles sur $X_\eta$) $= 1$ » : son $d'$ est ici l'indice, que le tapuscrit note $d^t$, et non l'ordre $d'$ de $R$ des pages 60 et 61.} Comme $\mathrm{Br}(K)$ est de
$p$-torsion, $0 \to \operatorname{Pic}(X_\eta) \to P(K) \to \Theta \to 0$ avec
$\Theta \subset \mathrm{Br}(K)$ de $p$-torsion, donc ${}_n\operatorname{Pic}(X_\eta)
\simeq {}_nP(K)$ pour $n$ premier à $p$ ; mais $\Gamma$ peut avoir de la torsion,
et l'on trouve la suite (6) de c), puis à la limite la suite (7) de d), avec $R =
\operatorname{Ker}(\Gamma \to \mathrm{NS}(X_0))$, et $R \otimes \mathbf{Z}_\ell =
\mathbf{Z}/\ell$ si $\ell \mid d$, $0$ sinon.\footnote{C'est l'encadré de la
marge. Il suppose $R$ égal au sous-groupe de torsion tout entier ; le tapuscrit,
plus prudent, écrit $\ell \mid d'$, $d'$ l'ordre de $R$.} On voit ainsi --- « au
moins modulo résolution des singularités », dit une incise --- que $R$ est de
torsion, donc un sous-groupe du $\mathbf{Z}/d$ de torsion de $\Gamma$, avec « une
tendance » à lui être égal. En marge : « demander à Raynaud si Ker est tout
$\mathbf{Z}/d\mathbf{Z}$ ».\footnote{En dimension relative 1, les pages 82 et 83
du même dossier répondent oui : leur $\Phi = \operatorname{Ker} u \simeq
\mathbf{Z}/d$ est ce noyau. Le rapprochement est nôtre.} La page esquisse enfin
la réduction au cas sans torsion par un scindage $\Gamma \simeq \Gamma' \oplus T$,
$T = \operatorname{Tors}\Gamma$, et s'arrête sur la suite $0 \to \delta' \to
\delta \to \delta'' \to 0$, $\delta'$ désignant la torsion première à $p$ de
$\mathrm{NS}(X_0)$.

\pagerange{72}{80}
\section*{Le théorème d'orthogonalité, démonstration arithmétique (pages 72 à 80)}

Une couverture : « Th.\ d'orthogonalité pour les modèles de Néron, démonstration
arithmétique », et plus bas, du même crayon, « $\ell$ (partiel !) ».\footnote{Un
signe au-dessus de « orthogonalité » se lit peut-être « $\ell/p$ ». Le
« (partiel !) » dit sans doute que seul le cas $\ell$ différent de la
caractéristique résiduelle est traité ; c'est ce que font les feuillets.} C'est
le théorème de l'exposé IX : pour une variété abélienne à mauvaise réduction, la
partie torique du module de Tate est orthogonale, pour l'accouplement de Weil, à
la partie fixe du module de la duale. La démonstration est « arithmétique » parce
qu'elle passe par les poids de Frobenius sur un corps fini.

\pagerange{73}{76}
\subsection*{L'énoncé général et sa démonstration (pages 73 à 76)}

\emph{Cadre.} Soit $A$ un anneau noethérien, $J$-adique, régulier et complet,
$X = \operatorname{Spec}A$, $Y = V(J)$, $U = X - Y$. Tout schéma fini étale sur
$Y$ se prolonge de façon unique en un schéma fini étale $\tilde Z$ sur $X$ ; de
même pour les faisceaux $\ell$-adiques lisses. Soient $G$, $G'$ des schémas en
groupes commutatifs plats sur $X$, dont les restrictions à $U$ sont des schémas
abéliens $A_U$, $A'_U$, munis d'une classe de correspondances divisorielles
$\xi$ sur $A_U \times A'_U$, et $\ell$ inversible sur $Y$, donc sur $X$. Soient
$H \subset G_Y$ un sous-groupe plat sur $Y$ dont les ${}_{\ell^\nu}H$ sont finis
sur $Y$, et $T' \subset G'_Y$ un tore.\footnote{La page note $S$ ce sous-groupe ;
on l'appelle $H$ pour garder $S$ aux bases. Le « NB » de la marge justifie la
suite : les ${}_{\ell^\nu}G$ sont étales et quasi-finis séparés sur $X$ et
contiennent donc les relèvements ${}_{\ell^\nu}\tilde H$.} Les modules de Tate
$T_\ell(\tilde H) \subset T_\ell(G)$ et $T_\ell(\tilde T') \subset T_\ell(G')$
sont définis, et leurs restrictions à $U$ sont des sous-modules de $T_\ell(A_U)$
et $T_\ell(A'_U)$. La classe $\xi$ définit l'accouplement de Weil
\[ \varphi_\xi\colon T_\ell(A_U) \times T_\ell(A'_U) \to
T_\ell(\mathbb{G}_{m,U}) . \]

\emph{Théorème.} $\varphi_\xi\bigl(T_\ell(\tilde H)|_U,\ T_\ell(\tilde
T')|_U\bigr) = 0$, en chaque point de $U$.

\emph{Démonstration.} Elle se fait en deux temps.

\emph{Réduction à un trait à corps résiduel fini.} Si $A/J$ est de type fini sur
$\mathbf{Z}$, tout point $x$ de $U$ a dans son adhérence un point $y$ de $Y$ de
corps résiduel fini ; on remplace $X$ par le complété en $y$, puis, $X - Y$ étant
de Jacobson, par l'adhérence normalisée d'un point $z$ avec
$\dim\overline{\{z\}} = 1$ : $X$ devient le spectre d'un anneau de valuation
discrète complet $V$ à corps résiduel fini $k$, et $Y$ son point fermé. Les
modules $T_\ell(\tilde H)$, $T_\ell(\tilde T')$ sont non ramifiés, et
l'accouplement, compatible à Galois, provient d'un accouplement
$\operatorname{Gal}(\bar k/k)$-équivariant
\[ T_\ell(H) \times T_\ell(T') \to T_\ell(\mathbb{G}_{m,k}). \]\footnote{La fin de la
page 73 n'est lue que par bribes ; les étapes de la réduction sont celles que la
page 74 rend lisibles.}

\emph{Lemme.} Soit $k$ un corps fini, $\ell$ différent de sa caractéristique,
$H$, $T$ deux groupes algébriques commutatifs sur $k$, $T$ de type multiplicatif.
Tout accouplement \emph{compatible à l'action de Galois}
\[ T_\ell(H) \times T_\ell(T) \to T_\ell(\mathbb{G}_m) \]
est nul.\footnote{La page dit « tout accouplement » ; c'est faux sans
l'équivariance, qui est l'hypothèse dont la démonstration se sert et que la
situation fournit.} En effet, on peut supposer $H$ et $T$ connexes et réduits, $T$
un tore déployé $\mathbb{G}_m^s$ après extension finie, et $H$ extension d'une
variété abélienne par un tore $\mathbb{G}_m^r$ (la partie unipotente n'a pas de
module de Tate en $\ell$). Le Frobenius agit par $q$ sur $T_\ell(\mathbb{G}_m^r)$,
sur $T_\ell(T)$ et sur $T_\ell(\mathbb{G}_m)$ ; un accouplement équivariant entre
les deux premiers vit sur un espace où Frobenius agit par $q^2$ et doit aboutir
là où il agit par $q$ : il est nul. Sur la partie abélienne, les valeurs propres
sont, par Weil, de valeur absolue $q^{1/2}$, donc $q^{3/2}$ sur le produit
tensoriel, et l'accouplement est encore nul. C'est un argument de poids au sens
d'aujourd'hui : $-1$ et $-2$ ne s'accouplent pas en $-2$.

\emph{Cas général.} On écrit $V = \varinjlim A_i$, $A_i \subset V$ de type fini sur
$\mathbf{Z}$, et $X = \varprojlim X_i$, $Y = \varprojlim Y_i$, $U = \varprojlim
U_i$. Pour $i$ grand, $G$, $G'$, $\xi$, $T'$ et $H$ proviennent de données sur
$X_i$ ; on peut supposer $H$ connexe (on le remplace par $H^0$), $\ell$ inversible
sur les $X_i$ (on remplace $A_i$ par $A_i[1/\ell]$), et l'on complète $A_0$ en
l'idéal de $Y_0$. On est ramené au théorème pour un anneau $\mathfrak{p}$-adique
régulier complet avec $A_0/\mathfrak{p}$ de type fini sur $\mathbf{Z}$, qui est le
premier cas. « CQFD ».\footnote{La page 76 est d'une lecture incertaine sur
plusieurs mots de liaison, et un « groupe » illisible suit « extension successive
d'une V.A., d'un tore et d'un » ; la structure de l'argument est claire.}

\pagerange{78}{80}
\subsection*{Les modèles de Néron (pages 78 à 80)}

\emph{Corollaire 1.} Soit $V$ un anneau de valuation discrète hensélien, de corps
résiduel $k$ et de corps des fractions $K$, $A$ une variété abélienne sur $K$, $G$
son modèle de Néron, $T$ le tore maximal de $G_k$, $\ell \neq \operatorname{car}
k$. On pose $M(A) = T_\ell(A(\bar K))$ ; on a
\[ M(A)^I_t \subset M(A)^I \subset M(A) , \]
où $M(A)^I = T_\ell(G_k)$ est la partie fixe sous l'inertie et $M(A)^I_t =
T_\ell(T)$ la partie torique.\footnote{Il note $M(A)^I_0$ la partie torique ; on
écrit $M^I_t$ pour ne pas suggérer une composante neutre. Le premier « $M(T)$ »
est corrigé en $M(A)$.} Si $A'$ est une variété abélienne sur $K$ et $\xi$ une
correspondance divisorielle sur $A \times A'$, d'où un accouplement compatible à
$\pi = \operatorname{Gal}(\bar K/K)$
\[ \varphi_\xi\colon M(A) \times M(A') \to T_\ell(\bar K^*) , \]
alors
\[ \varphi_\xi\bigl(M(A)^I,\ M(A')^I_t\bigr) = 0 \quad\text{et}\quad
\varphi_\xi\bigl(M(A)^I_t,\ M(A')^I\bigr) = 0 . \]
\footnote{La page annonce des « relations entre » ces modules, sans les écrire ;
ce sont celles-ci, que le théorème donne et que la page 80 énonce. Sous sa forme
publiée (SGA~7~IX), le théorème d'orthogonalité dit davantage, à notre
connaissance : lorsque $A'$ est la variété duale et $\xi$ la classe de Poincaré, $M(A)^I_t$ est exactement l'orthogonal de $M(A')^I$, et symétriquement.
Ces feuillets n'établissent que l'annulation.}

\emph{Corollaire 2.} Si $\xi$ est une polarisation de $A$, d'où une forme
bilinéaire alternée non dégénérée $\varphi_\xi$ sur $M(A)$, alors
$\varphi_\xi(M^I, M^I_t) = 0$.

\emph{Remarque.} Si $k$ est séparablement clos, les accouplements de niveau fini
\[ {}_{\ell^\nu}A'(K) \times {}_{\ell^\nu}T(k) \to \mu_{\ell^\nu} \]
ne sont pas nuls en général, même pour $\nu = 1$ : pour $x \in {}_{\ell^\nu}T(k)$
non nul, il existe $x' \in {}_{\ell^\nu}A'(\bar K)$ avec $\varphi(x, x') \neq 0$,
et, quitte à agrandir $K$, $x'$ devient rationnel. L'orthogonalité est un énoncé
$\ell$-adique, qui ne descend pas aux niveaux finis.

La page 80, barrée, reprend l'énoncé avec $V$ strictement hensélien, $A'$ la duale
et $G'$ son modèle de Néron, et la chaîne
\[ T_\ell(A(\bar K))^I = T_\ell(A(K)) = T_\ell(G(V)) = T_\ell(G(k)) =
T_\ell(G^0(k)) \supset T_\ell(T(k)) , \]
qui identifie la partie fixe au module de Tate de la fibre spéciale du modèle de
Néron ; puis elle recommence l'argument par passage à la limite, et s'arrête.

\pagerange{81}{84}
\section*{Autodualité de la jacobienne, suivant Artin--Mazur (pages 81 à 84)}

Une couverture porte, biffé, « Extensions et biextensions de groupes
algébriques \ldots », avec un mot encadré, « Obstruct\ldots », et, en dessous, le
titre retenu.\footnote{« Obstruct » est d'une lecture incertaine. Le dossier ne dit
pas de quel travail d'Artin et Mazur il s'agit.} Les deux pages qui suivent
calculent le groupe des composantes du modèle de Néron d'une jacobienne et la
dualité qu'il porte.

\pagerange{82}{82}
\subsection*{Le cas de la dimension relative 1 (page 82)}

Soit $X$ une courbe propre et plate sur un trait $S$, à fibre générique lisse et
géométriquement connexe, cohomologiquement plate en dimension 0, c'est-à-dire
$H^0(X_s, \mathcal{O}_{X_s}) = k(s)$.\footnote{« cohom.\ plat » est d'une lecture
incertaine ; la condition qui suit le précise. La lissité de la fibre générique
n'est pas écrite, et elle est nécessaire pour que $P^0_\eta$ soit une jacobienne.}
Soit $P = \underline{\operatorname{Pic}}_{X/S}$, $C_1, \ldots, C_r$ les
composantes de la fibre spéciale, de multiplicités $d_i$, et $\delta_i$ l'entier
tel que les degrés des faisceaux inversibles sur $C_i$ soient les multiples de
$\delta_i$. Soient $E \subset P$ le sous-groupe engendré par les classes des
$\mathcal{O}_X(C_i)$, $Q = P/E$, et
\[ P' = \operatorname{Ker}(P \xrightarrow{\deg} \mathbf{Z}_S), \qquad Q' =
\operatorname{Ker}(Q \xrightarrow{\deg} \mathbf{Z}_S) . \]
\footnote{La page ne définit ni $E$, ni $Q$, ni les $\delta_i$ ; ce sont ceux de la théorie de
Raynaud, et le N.B.\ de la page confirme $E$ et la division des degrés par les $\delta_i$.} On a un diagramme commutatif de
suites exactes de faisceaux fppf :

\begin{tikzcd}[column sep=small, row sep=small]
 & 0 \arrow[d] & 0 \arrow[d] & 0 \arrow[d] & \\
0 \arrow[r] & \Phi \arrow[r] \arrow[d] & P^0 \arrow[r] \arrow[d] & Q^0 \arrow[r] \arrow[d] & 0 \\
0 \arrow[r] & E \arrow[r] \arrow[d] & P' \arrow[r] \arrow[d] & Q' \arrow[r] \arrow[d] & 0 \\
0 \arrow[r] & E/\Phi \arrow[r] \arrow[d] & E^* \arrow[r] \arrow[d] & \Psi \arrow[r] \arrow[d] & 0 \\
 & 0 & 0 & 0 &
\end{tikzcd}

$E$, $E^*$, $\Phi$, $\Psi$ sont représentables par des schémas en groupes étales,
de fibre générique triviale. Si $k(s)$ est séparablement clos, leurs fibres
spéciales se calculent ainsi. Soient

\begin{tikzcd}[column sep=large]
0 \arrow[r] & \mathbf{Z} \arrow[r, "\lambda"] & \mathbf{Z}^r \arrow[r, "\alpha"] & E_{(s)} \arrow[r] \arrow[d, dashed, "u"] & 0 \\
 & \mathbf{Z} & \mathbf{Z}^r \arrow[l, "\mu"'] & E^*_{(s)} \arrow[l, "\beta"'] & 0 \arrow[l]
\end{tikzcd}

avec $\lambda(1) = (d_i)_i$, $\mu(n) = \sum_i d_i\delta_in_i$, $\alpha(n) =
\mathrm{cl}\,\mathcal{O}_X(\sum n_iC_i)$ et $\beta(L) = (\deg L|_{C_i}/\delta_i)_i$
; l'homomorphisme $u$ est donné par la matrice d'intersection :
\[ (\beta u\alpha)(n) = \Bigl(\sum_i n_i\,(C_i \cdot C_j)/\delta_j\Bigr)_{1
\leqslant j \leqslant r} . \]
Alors
\[ \Phi_{(s)} = \operatorname{Ker} u \simeq \mathbf{Z}/d, \qquad \Psi_{(s)} =
\operatorname{Coker} u, \qquad d = \operatorname{pgcd}(d_i) . \]
\footnote{Ce sont les identités qu'on vérifie : $\mu\beta u\alpha = 0$ parce que
$C_i \cdot X_s = 0$, et le noyau de la matrice d'intersection sur $\mathbf{Z}^r$
est $\mathbf{Z}(d_i/d)$, dont l'image dans $E_{(s)} = \mathbf{Z}^r/\mathbf{Z}(d_i)$
est cyclique d'ordre $d$. La page écrit aussi $\mu = (d^t_i)$.}

Si les anneaux locaux de $X$, après passage au hensélisé strict de la base, sont
factoriels, $\Psi_s$ est fini, comme $\Phi_s$ ; si de plus $k(s)$ est parfait ou
$(d^t, p) = 1$, $Q'$ est un modèle de Néron de
\[ Q'_\eta = P'_\eta = P^0_\eta = \underline{\operatorname{Pic}}^0_{X_\eta/K} .
\]
\footnote{Les biffures montrent l'hypothèse « $k(s)$ parfait ou $(d^t, p) = 1$ »
déplacée de la finitude de $\Psi_s$ à la propriété de Néron. C'est le théorème de
Raynaud (1970) sur le modèle de Néron d'une jacobienne ; voir Bosch,
Lütkebohmert, Raynaud, \emph{Néron Models}, chap.~9.} En particulier, si $d^t =
1$, donc $d = 1$ et $\Phi = 0$, $Q'$ est un modèle de Néron de $P^0_\eta$, pourvu
que $P'$ soit représentable.\footnote{La parenthèse « (loc.\ \ldots) » devant
« sur $S$ représentable » a un mot illisible, et « pourvu que » est d'une lecture
incertaine.} Le groupe $\Psi_s$ est donc le \emph{groupe des composantes} de la
fibre spéciale du modèle de Néron.

\pagerange{83}{84}
\subsection*{La dualité sur le groupe des composantes (page 83)}

Accouplons $\mathbf{Z}^r$ à lui-même par $\langle n, n'\rangle = \sum_i
\delta_in_in'_i$. Alors $\lambda$ et $\mu$ sont transposées, et $E^*(s) =
\operatorname{Ker}\mu$ s'accouple à $E(s) = \operatorname{Coker}\lambda$ par
$\langle\alpha(x), y^*\rangle = \langle x, \beta y^*\rangle$, à valeurs dans
$\mathbf{Z}$. De plus $u$ est \emph{autoadjoint} :
\[ \langle x, u(y)\rangle = \langle y, u(x)\rangle \qquad (x, y \in E(s)), \]
car pour $x = \alpha(n)$, $y = \alpha(n')$, $\langle n, \beta u\alpha(n')\rangle
= \sum_{i,j} n_jn'_i\,(C_i \cdot C_j)$, symétrique par symétrie de la forme
d'intersection.

Si les $\delta_i$ valent tous $1$ (par exemple $k(s)$ algébriquement clos), cet
accouplement identifie $E^*(s)$ à $\operatorname{Hom}(E(s), \mathbf{Z})$. Si de
plus $d = 1$, c'est-à-dire $d^t = 1$, $E(s)$ est libre, $u$ est injectif et
autoadjoint, et son conoyau $\Psi(s)$ hérite d'une \emph{autodualité canonique},
parfaite et symétrique, à valeurs dans $\mathbf{Q}/\mathbf{Z}$ : celle du groupe
discriminant d'une forme entière non dégénérée.\footnote{La page écrit « (p.ex.\
$k(s)$ parfait) » : pour un corps résiduel séparablement clos, parfait veut dire
algébriquement clos. La phrase qui ajoute $d = 1$ est d'une lecture incertaine
par endroits (« symétrique », un mot illisible).} « Cette autodualité exprime
l'obstruction à étendre la biextension qu'on a en $Q'_\eta \times Q'_\eta =
P^0_\eta \times P^0_\eta$ par $\mathbb{G}_m$, en une biextension des modèles de
Néron $Q'$, $Q'$ par $\mathbb{G}_m$. »

C'est, sous sa forme générale, l'accouplement que l'exposé IX définit sur les
groupes de composantes des modèles de Néron d'une variété abélienne et de sa
duale, comme obstruction à prolonger la biextension de Poincaré --- ce qu'on
appelle aujourd'hui \emph{l'accouplement de Grothendieck}. Ces deux pages en
donnent, pour une jacobienne, la description par la matrice
d'intersection.\footnote{Que l'accouplement de Grothendieck sur le groupe des
composantes d'une jacobienne coïncide avec celui que donne la matrice
d'intersection a été démontré par Bosch et Lorenzini (\emph{Invent.\ Math.} 148,
2002). La page affirme cette coïncidence sans la démontrer ; l'attribution « suivant
Artin--Mazur » de la couverture est la sienne.} La page 84 est une feuille
blanche à en-tête de l'IHES.

\end{document}
